\section{Introduction}

In this chapter, we describe our solution methodology for
LRSP. Although it would be possible to develop a branch-and-bound 
or a branch-and-cut algorithm for G-LRSP we suspect that such an 
approach would not yield positive results, given the inherent 
symmetry in the formulation, the weak LP relaxation bound and 
the known difficulty associated with solving related problems.
Instead, we use the set-partitioning-based formulation which is 
similar to those that have been successful for other hard 
combinatorial problems. 

As we have shown that adding both constraints (\ref{lb-open-fac}) 
and (\ref{relation}) to SPP-LRSP strengthen the formulation, 
to develop our solution methodology, we use an augmented SPP-LRSP 
formulation, denoted AUG-SPP-LRSP, which is comprised of the 
objective function (\ref{spp-obj}) and constraints (\ref{spp_con1}), 
(\ref{spp_con2}), (\ref{binary1}), (\ref{binary2}), (\ref{lb-open-fac}), 
and (\ref{relation}). 

%We have developed a branch-and-price algorithm based on formulation 
%AUG-SPP-LRSP to solve LRSP. 
In AUG-SPP-LRSP, each column corresponds to a feasible pairing.  
For medium- and large-sized instances, the number of feasible pairings 
can be extremely large, making it unlikely that we can efficiently solve 
instances that explicitly include all feasible pairings. Instead, we have 
developed a branch-and-price algorithm in which we dynamically generate 
only a subset of the feasible pairings at each node of the tree.


\section{Description of Solution Methodology} 
Branch-and-price that combine column generation with branch-and-bound 
algorithms has become a standard approach for solving large-scale 
integer programming  problems, especially those formulated using 
set-partitioning-based models. 
\citet{Barnhart} discuss the general methodology and the many  
challenges in developing a branch-and-price algorithm. The algorithm has been 
applied to a variety of problems, including various vehicle routing problems 
(\citet{Fukasawa}), location routing problems (\citet{Berger07}) and  
airline crew scheduling problems (\citet{Vance}), all of which are relevant in the 
context of the LRSP.

\zaupdate{Need more details about the literature?}
%\citet{Savelsbergh} applies a branch-and-price algorithm to a generalized 
%assignment problem which is a structure that appear in many models, such as vehicle routing, 
%facility location and scheduling. He showed that the algorithm outperforms many of the other 
%optimization algorithms used before. \citet{Vance} formulate an airline crew scheduling 
%problem and describe many details of their implementation of the branch and price algorithm. 
%Although the LP relaxation of their formulation gives a stronger bound than traditional 
%formulations, their LP relaxation is more difficult to solve than the traditional ones. 
%Closely related to our work is the work of \citet{Berger}, in which the authors 
%develop a branch and price algorithm for an uncapacitated location and routing problem.  

%[Although branch-and-price has been applied successfully to some of the 
%large scale optimization problems such as crew scheduling, machine
%scheduling, facility location and routing, since the mid-1990s, there are 
%many difficulties associated with the algorithm and the success of the 
%algorithm highly depends on the structure of the problem. \citet{Barnhart}
%provide an overview of the general methodology and many of the challenges.]

The basic idea of a branch-and-price algorithm is to solve an appropriate LP at 
each node of the branch-and-bound tree using column generation. 
Column generation is an approach for solving LPs for which we want to avoid 
explicitly enumerating all columns. The basic idea is to determine a subset of 
the variables that includes those that are nonzero in some optimal solution. 
By solving the LP including only this subset of columns rather than all possible 
columns, we can obtain the optimal LP solution. 

Column generation algorithms have been applied to a wide variety of large scale 
optimization problems, such as vehicle routing, crew scheduling, cutting stock, 
lot sizing and machine scheduling.  
%(\citet{Agarwal}, \citet{Desrochers2}), 
%airline crew scheduling (\citet{Anbil}, \citet{Vance}), location and routing (\citet{Berger05}), 
%cutting stock (\citet{Vance2}), lot sizing and machine scheduling (\citet{Vanderbeck}). 
\citet{Lubbecke} list other problems to which column generation has been applied.

Here, we explain the terminology and the main steps of the column generation algorithm.
In each iteration of the column generation procedure, we work with a restricted version 
of the LP relaxation; this {\em restricted master problem (RMP)} includes only a subset 
of the variables. We solve the RMP and use the associated {\em dual variable information} 
to formulate the {\em pricing problem} or {\em subproblem}. The objective of the pricing 
problem is to either identify new columns to add to the RMP or prove optimality of 
the current solution. For given dual variables, the objective function of the pricing 
problem is equal to the reduced cost function of a variable that we want to generate. 
The constraints of the pricing problem define the set of feasible columns corresponds 
to variables. From the theory of linear programming, we know that if every 
variable has a non-negative reduced cost, then the solution is optimal.
If the pricing problem identifies at least one column with negative cost, then we add
one or more columns to the RMP and reoptimize. Otherwise, the current solution is 
optimal for the full problem.

A column generation algorithm (COLGEN) for solving the LP relaxation of a formulation 
includes the following four steps.

{\bf COLGEN:}
\begin{enumerate}
\item[0.] Generate an initial set of columns that contains a feasible solution. Build a RMP.
\item[1.] Solve the RMP and obtain dual variable information.
\item[2.] Using the dual variables, form a pricing problem. 
  %that seeks to identify new pairings (columns) and their associated reduced costs.
\item[3.] Solve the pricing problem. 
  \begin{itemize}
  \item[i.]If the solution specifies a column with negative
    reduced cost, add the new column to the RMP and go to step 1. 
  \item[ii.]Otherwise, the solution found in step 1 is an optimal solution for 
    the LP relaxation. Stop.
  \end{itemize}
\end {enumerate}

\zaupdate{Mention challenges of COLGEN?}
One challenge of applying COLGEN is to solve the pricing problem which is generally
an integer programming model. One call to algorithm COLGEN may require many 
calls to the oracle that solves pricing problem. Therefore, an efficient solution method 
for pricing problem is crucial for the efficiency of the whole algorithm.   
In order to improve the performance of COLGEN,
one may prefer solving relaxations of the pricing problem which are easier to solve. 
However, replacing real pricing problem with a relaxed model yields a solution value
which is a lower approximation (for minimization problems) of the real LP relaxation 
solution. Another approach is to solve pricing problem heuristically. Indeed, step 3
of COLGEN does not require to find a column with  minimum cost (i.e., optimal solution 
to pricing problem) in intermediate iterations. However, to terminate the COLGEN 
in step 3 (ii) and conclude that the LP relaxation is optimal, one must find
a column with minimum cost at final iteration of COLGEN, thus one must solve
the pricing problem to optimality at least once.        
Another challenge about COLGEN is that executing step 1 may become time consuming
as the number of columns in RMP increases and column management may become a crucial part
of the algorithm. In column generation algorithms, it is also very common 
to iterate between steps 1 and 3 for a large number of iterations before termination. 
When there are alternative dual optimal solutions, the algorithm
jumps from one primal LP solution to another without an improvement in the
solution value.  

After describing column generation algorithm, we can list the steps of a 
branch-and-price algorithm which are similar to steps of a branch-and-bound.
The difference is that columns must be generated using COLGEN throughout the 
tree because we might not have all possible columns in the current node and branching 
rules and fixing of variables may change the values of the dual variables. 
%Therefore, the reduced costs of the columns may change 
%and there may exist new columns with negative 
%reduced costs that are not in the current problem. 
We use notation RMP$_2$ = COLGEN(RMP$_1$) to state that RMP$_2$ is
obtained after applying COLGEN to RMP$_1$. Figure 
\zaupdate{\ref{fig:bp_tree} shows a schematic representation of the  algorithm.
Update the figure}
{\bf BRANCH-PRICE}
\begin {enumerate}
\item[0.] Initialize a feasible RMP. Let the upper bound (UB) 
  be infinity and set of unprocessed nodes be empty.
  \item[]\noindent{\bf Root Node:} 
\item[1.] Call COLGEN(RMP) and let RMP =  COLGEN(RMP). 
  \begin{itemize}
  \item[i.] If the solution to RMP is integer, then STOP. The optimal
    integer solution is equal to RMP solution.
  \item[ii.] Otherwise, find a disjunction to eliminate the current fractional
    solution and create two new nodes corresponding to the disjunction. 
    Update set of unprocessed nodes with these two new nodes.
  \end{itemize}
\item [2.] Let lower bound equal to solution of RMP.
\item [3.] Call an oracle to find a primal feasible solution. If the
  solution is better than UB, update UB. 
 \item[] \noindent {\bf Node Selection \& Branching:}
\item [4.] If the set of unprocessed nodes is non empty, select one node
  to process, otherwise STOP. Return UB as an optimal integer solution.
 \item[] \noindent{\bf Non Root Node:} 
\item [5.] Modify RMP and the pricing problem in COLGEN according to the 
  disjunctions along the path in the tree from the root node to the 
  selected node. 
\item[6.]  Mark the node as processed. Call COLGEN(RMP) with modified 
  RMP and the pricing problem and let RMP =  COLGEN(RMP). 
  \begin{itemize}
  \item[i.] If the solution is infeasible, then fathom the node. Return to Step 4. 
  \item[ii.] If the solution to RMP is integer  then fathom the node. And if it is 
    better than UB, update UB.  Return to Step 4. 
  \item[iii.] If the solution value is greater than the UB, then fathom the node. 
    Return to Step 4. 
  \item [iv.] Otherwise, find a disjunction to 
    eliminate the current fractional  solution and create two new nodes 
    corresponding to disjunction. Update set of unprocessed 
    nodes with these two new nodes. Update lower bound based on
    the solution value.
  \end{itemize}
\item [7.] Call an oracle to find a primal feasible solution. If the
  solution is better than UB, update UB. Return to Step 4.  
\end{enumerate}

\begin{comment}   
\item {\bf Solving the Root Node of the Branch-and-Price Tree} 
  \begin {enumerate}
  \item [{\bf Step 1.}] Starting with an initial set of columns, solve the LP relaxation 
    of the AUG-SPP-LRS using a column generation algorithm.
  \item [{\bf Step 2.}] The LP relaxation solution is a lower bound for the problem.  
    If the LP relaxation solution satisfies the integrality constraints, 
    then it is an optimal solution of the original problem and the algorithm stops.  
    Otherwise, go to Step 3. 
  \item [{\bf Step 3.}] Find an integer feasible solution of the RMP with the current 
    set of columns using a branch-and-bound algorithm. Set this solution as the upper 
    bound for the problem. Go to step 4. 
  \end{enumerate}
\item {\bf Branching} 
  \begin{enumerate}
  \item [{\bf Step 4.}] Select one or more fractional variables on which to branch. 
    Create new nodes in the branch-and-price tree. 
  \item [{\bf Step 5.}] Do one of the following: 
    \begin {itemize}
    \item If there is an unprocessed leaf node in the tree, select one to process. Fix the 
      variables  according to the branching 
      rules applied along the path in the tree from the root node to the selected node. 
      Modify the pricing problem to incorporate the branching rules. Go to Step 6.
    \item If all nodes in the tree are marked as processed, then STOP. The current 
      best integer feasible solution is an optimal integer solution.
    \end{itemize}
  \item [{\bf Step 6.}] Solve the linear program associated with the current node using 
      a column generation algorithm with the modified pricing problem. 
      \begin {itemize}
      \item If the linear program at the current node is infeasible, then fathom the node,
	mark it as processed and return to Step 5. 
      \item If the LP solution at the current node satisfies the integrality constraints, 
        then it is a feasible solution for the original problem. If the current solution
        is better than the best integer feasible solution, then update the upper bound. 
	Return to Step 5. 
      \item If the LP solution contains fractional variables, go to Step 4. 
    \end{itemize}
  \end {enumerate}
\end{enumerate}
\end{comment}

\begin{figure}[h]
  \centering
  \includegraphics[scale=0.5]{figs/figure2.eps}
  %\fbox{ \includegraphics[scale=0.5]{figure2.eps}}
  \caption{Branch-and-Price Tree}\label{fig:bp_tree}
\end{figure}

Besides the challenges in column generation algorithm, the main challenge in developing a 
branch-and-price algorithm is developing effective branching rules that do not destroy the 
structure of the pricing problem. At non root nodes of the tree, the RMP is modified based
on the branching rules applied along the path and so is pricing problem. Therefore, 
applied branching rules must not complicate the pricing problem and can easily be 
incorporated in the solution algorithm for pricing problem. 
In the following sections, we describe the components of the branch-and-price algorithm
that we developed for the LRSP. 
\zaupdate{In Section \ref{sec:PricingProb}, we describe the column generation 
algorithm for solving the LPs at the nodes of the branch-and-price tree. 
In Section \ref{sec:BranchingRules}, we describe the branching rules. 
In Section \ref{sec:BPAlgorithm}, we present implementation details of the algorithm. 
}


\section{Column Generation}
To begin the column generation procedure, we need to construct an initial RMP
for AUG-SPP-LRSP that includes all of the facility location variables but
only a subset of the pairing variables. We solve the RMP and use the associated 
dual variable information to formulate the pricing problem. 
In the context of the AUG-SPP-LRSP, the objective of the pricing problem is to 
determine if there are feasible pairings (based on the definition and constraints 
in Section \ref{set-part-form}) with negative reduced cost that can be added to the
RMP. If the pricing problem identifies at least one pairing with negative 
reduced cost, then we add one or more columns to the RMP and reoptimize.
Otherwise, the current solution is the optimal solution for the full model.

\subsection {Formulating the Pricing Problem}
\label{form-pricing-prob}
As explained in Section \ref{set-part-form}, set of pairings for each facility
are independent, thus we decompose the pricing problem into a set of independent 
pricing problems, one for each facility. To write the objective function of pricing
problem for facility $j$, we consider formulation AUG-SPP-LRSP and obtain an 
expression for the reduced cost of a pairing. Let $\pi_{i}$ be the dual variable 
associated with constraint (\ref{spp_con1}) for customer $i$, $\mu_{j}$ be 
the dual variable associated with constraint (\ref{spp_con2}) for facility $j$, 
and $\sigma_{ji}$ be the dual variable for linking constraint (\ref{relation}) for 
facility $j$ and customer $i$. For the variable $z_{p}$ for $p \in P_{j}$, 
the reduced cost $\hat{C}_{p}$ can be written as   
\begin{eqnarray}
\label{red-cost}
\hat{C}_{p}=C_{p}-\sum_{i\in I}a_{ip} \pi_{i}+\sum_{i \in I}a_{ip} D_{i} \mu_{j}+\sum_{i \in I} a_{ip} \sigma_{ji},
\end{eqnarray} 
where $C_{p}$ is the cost of pairing $p$. Then, we can express the pricing problem 
for facility $j$ as the problem of minimizing $\hat{C}_p$ subject to 
the following constraints defining feasible pairings. 
\begin{itemize}
\item [i.] Each route included in the pairing starts and ends at the 
same facility, 
\item [ii.]for each route in the pairing, the total demand of the 
customers assigned to the route is at most the vehicle capacity, 
\item [iii.] the total travel time of the pairing is at most the vehicle time limit,
\item [iv.] each customer included in the pairing must be visited only once. 
\end{itemize}
We can formulate this problem for facility $j$ as follows:
\begin{optprog}
 {\bf SP} Minimize & \objective{C^V + \sum_{i \in M}\sum_{k \in M} \hat{c}_{ik}x_{ik}}\label{pricing-j-0}\\
  &\sum_{k \in M}x_{ik} & \leq & 1 & \forall i \in I \label{pricing-j-1}\\
  & \sum_{k\in M} x_{ik}-\sum_{k \in M} x_{ki} & = & 0 & \forall i \in M, \label{pricing-j-2}\\
  & y_{ik} - L^V x_{ik} & \leq & 0 & \forall i \in M, k \in M, \label{pricing-j-3}\\
  & \sum_{k \in M} y_{ik} - \sum_{k \in M} y_{ki} + D_i\sum_{k \in M}x_{ik} & = & 0 & \forall i \in I, \label{pricing-j-4}\\
  & \sum_{i \in M}\sum_{k \in M} T_{ik} x_{ik} & \leq & L^T, &  \label{pricing-j-5}\\
  & x_{ik} & \in & \{0,1\} & \forall i \in M, k \in M, \label{pricing-j-6}\\
  & y_{ik} & \geq &  0 & \forall i \in M, k \in M, \label{pricing-j-7}         
\end{optprog}
where $M = I \cup \{j\}$, $x_{ik} = 1$ if link $(i,k)$ is used in solution, 
and 0 otherwise, $y_{ik}$ is equal to the amount of flow on link $(i,k)$,
and $\hat{c}_{ik}$ is the cost of traveling link $(i,k)$. 
We refer  $\hat{c}_{ik}$ as the {\em reduced cost of link $(i,k)$} and calculate
as follows based on the reduced cost function \ref{red-cost}. 
\begin{eqnarray}
\hat{c}_{ik} & = & \left \{  \begin{array}{ll}
                 C^O T_{ik}-\pi_{k} + D_k \mu_j + \sigma_{jk} & \mbox{if $i \in M$, and $k \in I$}, \\
                 C^OT_{ij} & \mbox{if $i \in I$}.
                 \end{array}
              \right. \label{pricing-j-0-open}
\end{eqnarray}
To write the objective function (\ref{pricing-j-0}), recall that in equation 
(\ref{red-cost}), the cost $C_{p}$ is the sum of the fixed cost of a vehicle 
and the variable costs of the routes that are included in the pairing. The variable 
cost of a route is equal to the sum of the costs of arcs in the route. 
We can include part
$-\sum_{i\in I}a_{ip} \pi_{i}+\sum_{i \in I}a_{ip} D_{i} \mu_{j}+\sum_{i \in I} a_{ip} \sigma_{ji}$
to the objective function by transferring the values $\pi_i$, $D_{i} \mu_{j}$ and $\sigma_{ji}$ 
associated with customer $i$ as costs to the incoming arcs of customer node $i$. 
Therefore, in the definition of $\hat{c}_{ik}$ in (\ref{pricing-j-0-open}),
operating cost of arcs into customer nodes are modified with corresponding dual variables.
With this modification, the reduced cost of a pairing can be computed as the 
sum of the reduced costs of the arcs in the pairing.     


\subsection{Elementary Shortest Path Problem With Resource Constraints (ESPPRC)}
The pricing problem (SP) is a mixed integer programming model. On the other hand,
it  can be defined as an elementary shortest path problem with resource 
constraints (ESPPRC) on a modified network. In an ESPPRC, the objective is to find a 
shortest path between the source node and the sink node such that each other node 
is visited at most once and the resource constraints are not violated. 

To properly represent our pricing problem as an ESPPRC, we need to make one adjustment; 
while in an ordinary elementary shortest path problem, each node can be visited at most 
once, in our case, because we are looking for a set of routes, we allow the sink node to 
be visited more than once. In our modified ESPPRC, the sink node represents the end of 
one route and maybe the beginning of another. In an ordinary ESPPRC, a path from 
the source to the sink corresponds to one route, however, a solution of our modified 
ESPPRC may correspond to multiple routes. Another reason for multiple visits to the sink 
is to reset the remaining capacity of the vehicle before the new route. In a pairing, 
the vehicle starts a route with zero truck load and picks up the required demand  
from all of the customers on the route. The vehicle returns to the depot after the last 
customer in the route and resets its capacity to zero load for the next route. 
In the pricing problem, the truck capacity is one of the resources. Whenever the 
path visits the sink in the modified network, the vehicle capacity resource consumption 
is updated and set to zero load. In the next section, we describe a network
for ESPPRC to solve pricing problem (SP).

\subsection{Underlying Network Structure For Modified ESPPRC}
\label{network-str-for-1-stp}
To define the pricing problem for facility $j$, we construct a network with $|I|+2$ nodes, 
$|I|$ nodes for the customers, one for facility $j$ as a source node and a copy of facility 
$j$ as a sink node. Let $A^j$ be the arc set of the network for facility $j$, $s$ represent 
the source node, $si$ represent the sink node. Then $A^j$ consists of four sets of arcs: 
%\begin{eqnarray*}
\begin{optprog}
  & A^{j}_{s,c}  & = & (\{s\} \times I)   & = \mbox{set of arcs from the source to customer nodes} \\ 
  & A^{j}_{c,si} & = & (I \times \{si\})  & = \mbox{set of arcs from customer nodes to the sink} \\
  & A^{j}_{si,c} & = & (\{si\} \times I)   & = \mbox{set of arcs from the sink to customer nodes} \\
  & A^{j}_{c,c}  & = & (I \times I)    & = \mbox{set of arcs between pairs of customer nodes}
%\end{eqnarray*}
\end{optprog}
Note that there are no arcs from customer nodes back to the source but there are arcs from 
the sink to customer nodes, this enables us to generate pairings including multiple routes.
In a path from source in this network, the arcs from the sink to customer nodes represent 
the beginning of a new route, whereas the arcs from the customer nodes to the sink represent 
the end of a route. We let each customer node has a demand equal to its demand in the LRSP 
instance. Demand of the source is zero while demand of the sink is equal to minus vehicle 
capacity $-L^V$. Vehicle capacity is a resource constraint in the problem and total demand 
must be less than the vehicle capacity  on a path from source to any node in the network. 
And in the real problem, total demand of each route must be less than the vehicle capacity. 
Since the arcs from the sink to customer nodes represent the beginning of a new route, 
the vehicle load, total demand at the beginning must be reset to zero at sink node. 
By setting total demand to maximum of zero and total demand at any node, and letting 
the demand of sink node be minus vehicle capacity, a new route starts with zero 
total demand. 

To generate pairings that can be served within time limit of a vehicle, we let the distance 
between each pair of nodes be equal to the travel time between these nodes given in the 
LRSP instance. Time limit is the other resource constraint in the problem and total length
of any path in the network must not exceed the time limit. Finally, to obtain 
that the cost of a path from source to sink in the constructed network is equal to the 
reduced cost of the corresponding pairing, we use the following cost function based
on the objective function of (SP). 
\begin{eqnarray}
c^{'}_{ik} & = & \left \{  \begin{array}{ll}
                 c_{ik}-\pi_{k} + D_k \mu_j + \sigma_{jk} & \forall (i,k) \in  (A^{j}_{s,c} \cup A^{j}_{si,c} \cup A^{j}_{c,c}), \\
                 c_{ij} & \forall (i,k) \in A^{j}_{c,si},
                 \end{array}
              \right. \label{cost-espprc-ntw}
\end{eqnarray}
where $c_{ik} = C^O T_{ik}$ be the operating cost of link $(i,k)$ in LRSP instance.
%of minimum reduced cost using the constructed network, we need to find an 
%appropriate way to incorporate the dual variables into the costs of the arcs in  $A^j$. 
%Remember that $t_{ik}$ is the travel time from location $i$ to location $k$ and $d_i$ is the 
%demand of customer $i$ and let $c_{ik}$ be the operating cost of traveling from location 
%$i$ to location $k$ in the LRS problem. To define the length and the cost of arcs on the 
%constructed network, let $t^{'}_{ik}$ be the travel time on arc $(i,k)$ and $c^{'}_{ik}$ 
%be the reduced cost of traveling arc $(i,k)$. Then, for the network associated with 
%facility j:  
\begin{comment}
\begin {itemize}
\item {\em Case 1:} If $(i,k)$ is an arc from source or sink to a customer node, 
  e.g., $(i,k) \in (A^{j}_{s,c} \cup A^{j}_{si,c})$, then 
  $t^{'}_{ik}=t_{jk}$, $c^{'}_{ik}=c_{jk}-\pi_{k}+d_k \cdot \mu_j+\sigma_{jk}$ 
\item {\em Case 2:} If $(i,k)$ is an arc between two customer nodes, 
  e.g., $(i,k) \in A^{j}_{c,c}$, 
  then   $t^{'}_{ik}=t_{ik}$, $c^{'}_{ik}=c_{ik}-\pi_{k}+d_k \cdot \mu_j+\sigma_{jk}$ 
\item {\em Case 3:} If $(i,k)$ is an arc from a customer node to sink, 
  e.g., $(i,k) \in A^{j}_{c,si}$, then $t^{'}_{ik}=t_{ij}$, $c^{'}_{ik}=c_{ij}$
\end{itemize} 
We transfer the dual variables $\pi_i$ and $\sigma_{ji}$ associated with customer 
$i$ to the incoming arcs of customer node $i$. As in the reduced cost calculation
in (\ref{red-cost}), the dual variable for facility $j$, $\mu_j$, is multiplied with
demand value $d_i$ for customer $i$ that is in the pairing. Therefore, the cost 
$d_i \cdot \mu_j$ is also carried on to the incoming arcs of customer node $i$. Since facility
$j$ does not have demand, the dual variable $\mu_j$ does not affect the arcs 
into the sink node.  
\end{comment}

A path from the source to the sink on the constructed network corresponds to a pairing,
if (i) it visits any customer node at most once, (ii) total demand from source
to any node in the path does not exceed the vehicle capacity, and (iii) 
the total length of the path does not exceed the time limit. These conditions
are satisfied if the path is a feasible solution to ESPPRC on this network.
And total cost of the path plus the vehicle fixed cost $C^V$ equals to the reduced cost 
of the pairing for a given dual vector. If the path visits 
the sink  more than once, then it corresponds to a pairing that includes more than 
one route. Figure \ref{fig:exp_network} shows an example network and a path on the 
network. In the example, the path corresponds to pairing 
(facility$_j$ - 3 - facility$_j$ - 1 - 2 - facility$_j$) which includes two routes. 
\begin{figure}[h]
  \centering
  \includegraphics[scale=0.8]{./figs/NetworkexpG-2.eps}
  \caption{Example of a network and a path from source to sink (dashed line)} \label{fig:exp_network}
\end{figure}

%A visit to the sink resets the remaining capacity of the vehicle to a full truckload. 
%In a pairing, the vehicle starts a route with a full truck 
%load and delivers the required demand to all of the customers on the route. The vehicle returns 
%to the depot after the last customer in the route and resets its capacity to full truck load for 
%the next route. In the pricing problem, the truck capacity is one of the resources. Whenever the 
%path visits the sink in the modified network, the vehicle capacity resource consumption is updated and 
%set to full truck load. 
%$A^{j}_{si,c}$, arcs from sink to customers and 
%$A^{j}_{c,c}$, arcs between any two customers. There are no arcs from customer 
%nodes to the source. 

By using the example in Figure \ref{fig:exp_network} we can easily illustrate
that the total cost of the path is equal to the reduced cost of the corresponding pairing.
Using the cost function defined in (\ref{cost-espprc-ntw}), the total cost
of the path in the example plus the fixed cost is:
\begin{eqnarray}
\mbox{Total cost of the path} + C^V & = & c^{'}_{j,3}+c^{'}_{3,j}+c^{'}_{j,1}+c^{'}_{1,2}+c^{'}_{2,j} + C^V \nonumber\\
                              & = & (c_{j,3}-\pi_3+D_3\mu_j+\sigma_{j,3})+ c_{3,j} + (c_{j,1}-\pi_1+D_1\mu_j+\sigma_{j,1}) + \nonumber\\
                              & + & (c_{1, 2}-\pi_2+D_2\mu_j+\sigma_{j,2}) + c_{2,j} + C^V \label{exp-red-cost-clc}
\end{eqnarray} 
Grouping terms differently and noting that the column vector representing this pairing is 
$a_{p}=[1 1 1]$ and the cost of a pairing is the sum of the arc costs plus the vehicle 
fixed cost, we obtain the following expression, which is just the reduced cost of $z_{p}$ as calculated
via equation (\ref{red-cost}).
\begin{eqnarray*}
(\ref{exp-red-cost-clc})  & = & VF + (c_{j,3}+c_{3,j}+ c_{j,1}+ c_{1,2} + c_{2,j}) \\
                & - & (\pi_3+\pi_1+\pi_2)+ \mu_j (D_3+D_1+D_2)  +  (\sigma_{j,3}+\sigma_{j,1}+\sigma_{j,2})\\
             &=& C_{p}-\sum_{i\in I}a_{ip} \pi_{i}+\sum_{i \in I}a_{ip} D_{i}\mu_{j}+\sum_{i \in I}a_{ip}\sigma_{ji} = \hat{C}_{p}
\end{eqnarray*} 
Therefore, the reduced cost of a pairing in the original problem is equal to the cost
of a path in the modified network.

\begin{comment}
To illustrate why a path through the modified network is equivalent to a pairing, we consider
the following example. Let $N=\{1,2,3,4,5,6,7,8,9,10\}$ be the set of demand nodes.  Then the 
set of all nodes in the network is $I=\{source,1,2,3,4,5,6,7,8,9,10,sink\}$. Consider the path
(source-2-4-sink-5-6-7-sink-9-10-sink) in the modified network. This path in the modified network
is equivalent to a pairing with three routes in our pricing problem. The component routes are 
($facility_j$-2-4-$facility_j$), ($facility_j$-5-6-7-$facility_j$) and ($facility_j$-9-10-$facility_j$). 
The pairing is feasible if each route's total demand does not exceed the vehicle capacity and if the 
sum of the travel times of the paths does not exceed the time limit. From the arc costs  specified 
in Cases 1 through 3, the reduced cost of the path in the modified network is:
\begin{eqnarray*}
\mbox{Total cost of the path} &=& VF+(c^{'}_{j,2}+c^{'}_{2,4}+c^{'}_{4,j}+c^{'}_{j,5}+c^{'}_{5,6}+c^{'}_{6,7}+ c^{'}_{7,j}+c^{'}_{j,9}+c^{'}_{9,10}+c^{'}_{10,j})\\
                              &= &VF+(c_{j,2}-\pi_2+d_2\mu_j+\sigma_{j,2})+(c_{2,4}-\pi_4+d_4\mu_j+\sigma_{j,4}) \\
                              &+ &(c_{4,j})+(c_{j,5}-\pi_5+d_5\mu_j+\sigma_{j,5})+(c_{5,6}-\pi_6+d_6\mu_j+\sigma_{j,6})\\
                              &+ &(c_{6,7}-\pi_7+d_7\mu_j+\sigma_{j,7})+(c_{7,j})+(c_{j,9}-\pi_9+d_9\mu_j+\sigma_{j,9})\\
                              &+ &(c_{9,10}-\pi_{10}+d_{10}\mu_j+\sigma_{j,10})+(c_{10,j})
\end{eqnarray*} 
Grouping terms differently and noting that the column vector representing this pairing is 
$a_{jp}=[0 1 0 1 1 1 1 0 1 1]$ and the cost of a pairing is the sum of the arc costs plus the vehicle 
fixed cost, we obtain the following expression, which is just the reduced cost of $Z_{jp}$ as calculated
via equation (\ref{red-cost}).
\begin{eqnarray*}
\hat{C}_{jp} &=& VF + (c_{j,2}+c_{2,4}+c_{4,j}+c_{j,5}+c_{5,6}+c_{6,7}+ c_{7,j}+c_{j,9}+c_{9,10}+c_{10,j})\\
             &-& (\pi_2+\pi_4+\pi_5+\pi_6+\pi_7+\pi_9+\pi_{10})+\mu_j \cdot (d_2+d_4+d_5+d_6+d_7+d_9+d_{10})\\
             &+& (\sigma_{j,2}+\sigma_{j,4}+\sigma_{j,5}+\sigma_{j,6}+\sigma_{j,7}+\sigma_{j,9}+\sigma_{j,10})\\
             &=& C_{jp}-\sum_{i\in N}a_{ipj} \cdot \pi_{i}+\sum_{i \in N}a_{ipj} \cdot d_{i}.\mu_{j}+\sum_{i \in N}a_{ipj} \cdot \sigma_{ji}
\end{eqnarray*} 
\end{comment}


In the following section, we describe approaches for solving the pricing problem defined on the
constructed network. 

\section {Solving the Pricing Problem Exactly}
\label{solv-pri-exact}
The pricing problem is formulated as an ESPPRC. If the underlying network is acyclic, then
the ``elementary'' property of the paths can be ignored, since than solution of a shortest path 
problem with resource constraints (SPPRC) is equal to the solution of ESPPRC. 
\citet{Desrochers} developed a label correcting algorithm for SPPRC that is based on dynamic 
programming and that has pseudo polynomial complexity depending on size of the resource
windows and the number of resources. The algorithm has been shown to be  successful when there 
are tight resource constraints, but it becomes more time consuming as the number of resources 
increases. As the algorithm builds partial paths, it assigns a label to each 
indicating the resource consumption. To eliminate the problem of an exponentially 
increasing number of labels, the algorithm applies dominance rules to delete dominated labels. 

If the network is not acyclic, which is generally the case in pricing problems because of the 
negative arc costs, ``elementary'' property of the paths must be considered. \citet{Beasley} 
describe how to find elementary paths in a graph by adding an extra binary resource for each 
node showing whether the node is visited or not. However, this increases the number of resources
and possible combinations exponentially. In fact, ESPPRC is proved to be NP-Hard (\citet{Dror}).
\citet{Feillet} adapt \citeauthor{Desrochers}'s label correcting algorithm to find elementary paths 
and strengthen the algorithm by developing additional domination rules for the labels. 
We have adapted \citet{Feillet}'s label correcting algorithm to our pricing problem.
% to solve the  ESPPRC described in Section \ref{form-pricing-prob}. 

We have developed two solution approaches. One approach is to solve one ESPPRC on the network 
described in Section (\ref{network-str-for-1-stp}). An alternative approach is to extend the
first approach to a two-phase pricing problem in which  each  phase is also formulated as a 
network problem and solved as an ESPPRC. 

\subsection{1p-ESPPRC: One Phase Pricing Problem}
\label{1pespprc}
In this approach we extend \citet{Feillet} algorithm to solve one ESPPRC on the constructed 
network. In the \citet{Feillet} algorithm, each path from the source to a node in the network is
assigned a label, which is a vector of the reduced cost of that (partial) path, the resources 
consumed (vehicle load and elapsed time), the customer nodes included in the path, and the 
customer nodes that cannot be visited due to the resource constraints. Customers are 
considered {\em unreachable} if they have already been visited in the path or if they cannot 
be served by the vehicle without violating a resource constraint. The algorithm fans out from 
the source node, repeatedly examining the nodes of the graph. Each time it examines a node, 
the algorithm extends each of the paths to each possible successor node, continuing until no 
further extensions are feasible.

To explain the algorithm more specifically, we let $l_r^i$ be a label corresponds to 
path $r$ from source to node $i$. We state that vector 
$l_r^i = (R^i_c, R^i_{v},R^i_{t}, R^i_n, V^i_1,..,V^i_{|I|})$ keeps the information
about path $r$ where
\begin{eqnarray*}
  R^i_c    & = & \mbox{reduced cost of the path,} \\
  R^i_{v}  & = & \mbox{total demand of the path,}  \\
  R^i_{t}  & = & \mbox{time consumption of the path,}  \\
  R^i_n    & = & \mbox{total number of unreachable customer nodes,}  \\ 
  V^i_k & = & \left\{ \begin{array}{llll}
                  -2 & \mbox{if customer $k$ is temporarily unreachable for}\\
                     & \mbox{the path}\\
                  -1 & \mbox{if customer $k$ is unreachable for the path}  \\
                     & \mbox{because of the resource limits,} \\
                  0  & \mbox{if customer $k$ is not in the path and it is}  \\
                     & \mbox{reachable for the path with the resources,}  \\
                 s>0 & \mbox{if customer $k$ is the $s^{th}$ node in the path,} 
                \end{array} 
              \right\} \forall k \in I. 
\end{eqnarray*}
% $R^i_c$ is the reduced cost, $R^i_{v}$ is the current,R^i_{t}, R^i_n
%Let $X_{so,i}$ be a path from the source node to node $i$. The label $L^{X_{so,i}}$ 
%that corresponds to path $X_{so,i}$, is associated with a state $(T_i^{v},T_i^{t},s_i,V_i^1,..,V_i^N)$  
%that keeps the information about the path. 
%In the state,  $T_i^{v}$ is the current vehicle load,
%$T_i^{t}$ is the time consumption of the path and the vector $(V_i^1,..,V_i^N)$ gives
%the sequence of nodes in the path as well as the unreachable nodes for the path. 
%The elements of the vector, $(V_i^1,..,V_i^N)$ are defined as:
%\begin{eqnarray*}
%V_i^k & = & \left\{ \begin{array}{llll}
%                 -1 & \mbox{if node $k$ is {\bf unreachable} from node $i$ for the path because of the}\\
%                    & \mbox{ resource limits, }\\
%                  0 & \mbox{if node $k$ is not in the path and it is {\bf reachable} from node $i$ with the }\\
%                    & \mbox{current resources,}\\
%                 n>0 & \mbox{if node $k$ is {\bf unreachable} from node $i$ since it is in the path and it is} \\
%		     & \mbox{the $n^{th}$ node.} 
%                \end{array} 
%              \right. 
%\end{eqnarray*}
%The state variable $s_i$ keeps the number of unreachable nodes. In the label $L^{X_{so,i}}$,  node $k \in N$ is 
%{\em unreachable} from node $i$ if it is already in the path or if there is not 
%enough of a resource to extend the path (label) from node $i$ to node $k$. 
Note that the dimension of the vector $V$ is $|I|$, there is one member
for each customer node in the network. In label $l_r^i$, if customer $k$
is visited in the path, $V^i_k$ keeps the place of the customer in the path and
remarks that customer $k$ is unreachable for this label because of the elementary 
property. Based on \citet{Feillet}'s algorithm, if customer $k$ is not visited in 
the path $r$, $V^i_k$ is set to either 0 or -1. Even though we have two resource constraints 
(vehicle capacity and time limit) in the problem,  we only check the time limit constraint
to mark a customer as ``unreachable'', i.e., to set $V^i_k = -1$. 
Since if vehicle load exceeds vehicle capacity by adding customer $k$ to the path, 
that does not mean that we can not add customer $k$ to this path in the future. It 
may be possible that a path extended from path $r$ stops at sink node and would have
enough capacity to visit customer $k$ in further extensions. 
In 1p-ESPPRC, we extended \citet{Feillet}'s algorithm by adding a new domination rule
which will be explained later in this section. For this extension, if customer $k$
satisfy the conditions in the domination rule, we set $V^i_k$ to -2 meaning that 
the customer is {\em temporarily unreachable} and this condition may change further in
the path. We do not need an element in $V$ for the source, because every path starts 
from the source and can never return to it. Similarly, we do not need an element for 
the sink, because the sink must always be reachable for a path. Extending a path to 
the sink is always resource feasible: (1) since demand of the sink is minus 
vehicle capacity, thus total demand would be less than the vehicle capacity and 
(2) while constructing label $r$ to a non sink node $i$ 
($l_{r}^i = (R^i_c, R^i_{v},R^i_{t}, R^i_n, V)$), the time limit constraint is checked 
to ensure that the path is further extendable to the sink, i.e., 
$R^i_{t} + T_{i,si} \leq L^T$; otherwise it is not created since every path
must reach to the sink at the end.     

%a new path $r$ to a non sink node $i$, the time constraint 
%is checked to ensure that the path is further extendable to the sink, 
%i.e., $R^i_{t} + T_{i,si} \leq L^T$; otherwise it is not created since every path
%must reach to sink at the end.     

%For example, $k$ is unreachable
%if the vehicle load is not enough to serve the demand of customer $k$ 
%($T_i^v-d_k<0$) and it is not possible to 
%travel to node $k$ and then to the sink (to finish the path) within the time 
%limit ($T_i^{t}+t'_{ik}+t'_{k,sink}>TL$).   

The label correcting algorithm for solving the ESPPRC includes three main parts: \\
{\bf 1p-ESPPRC:}
\begin{enumerate}
\item {\bf Initialization step} 
  \begin{enumerate}
  \item To begin the algorithm, a label list is created 
    for each node to keep the generated labels (partial paths to the node). 
   % Let $L_k$ be the label list for node $k$, for all $k \in I \cup \{s,si\}$.
  \item %Set $L_k = \emptyset$ for all  $k \in I \cup \{si\}$,
    All label lists except the list for the source node are initially empty. 
    %and $L_{s} = \{l^s_1\}$ where $l^s_1 = 0^{(1 \times |I|+4)}$, i.e., 
    The source has one label which is a zero vector, i.e., resource consumptions 
    and the reduced cost  are zero, and all customers are reachable.
  \end{enumerate}
\item {\bf Extension step} 
  \begin{enumerate}
  \item For each node, the label list is searched and all unprocessed labels 
    are extended to all of the possible successors of the node. Unprocessed means 
    that the label was not extended in previous iterations. 
    The label $l^i_r$ for node $i$ can be extended to node $k$ if 
    %one of the following is true: 
    \begin{itemize}
    \item[i.] reachable: $k = si$ or $V^i_k = 0$, and
    \item[ii.] resource feasible: $R^i_{v} + D_k \leq L^V$ and 
      $R^i_{t} + T_{i,k} + T_{k,si} \leq L^T$.
      %$T_i^v-d_k\geq 0$ and $T_i^{t}+t'_{ik}+t'_{k,sink}\leq TL$
      % the node $k$ is reachable from node $i$ ($V_i^k=0$) and the extended label is resource 
      % feasible ($T_i^v-d_k\geq 0$ and $T_i^{t}+t'_{ik}+t'_{k,sink}\leq TL$). 
    \end{itemize}
  \item If it is feasible to extend $l^i_r$ to node $k$, a new label is
    created for node $k$ by updating the resource consumptions.
    Let $l^k_p$ be the label extended from  $l^i_r$, then
    \begin{eqnarray*}
      R^k_c    & = & R^i_c + c^{'}_{i,k}, \\
      R^k_{v}  & = & \mbox{max}\{R^i_v + D_k, 0 \},  \\
      R^k_{t}  & = &  R^i_t + T_{i,k},  \\
      V^k_g & = & \left\{ \begin{array}{llll}
                N & \mbox{ if $g = k$ and $N = \#$ of nodes in the path,}  \\
               -1 & \mbox{if  $V^i_g = 0$ and $R^k_{t} + T_{k,g} + T_{g,si} > L^T$,}  \\
            V^i_g & \mbox{if  $V^i_g \neq 0$,}\\
                0 & \mbox{otherwise,}
      \end{array} 
      \right\} \forall g \in I, \\
     R^k_n    & = &  \sum_{g \in I: V^i_g \neq 0} 1. 
    \end{eqnarray*}
%    $R^k_c = R^i_c + c^{'}_{i,k}$, $R^k_v = \mbox{max}\{R^i_v + D_k, 0 \}$,
%    $R^k_t = R^i_t + T_{i,k}$. If $k$ is a customer node then set 
%    $V^k_k$ to the number of nodes in the path $p$.  If $V^i_g = 0$ 
%    and $R^k_{t} + T_{k,g} + T_{g,si} > L^T$ for $g \in I$, then  
%    $g$ is unreachable for path $p$, set $V^k_g = -1$. Otherwise,
%    and $V^i_g = V^k_g$. 
  \item Mark the labels that are extended as processed.
  \end{enumerate}
\item {\bf Elimination step } 
  \begin{enumerate}
  \item 
    After extending unprocessed labels of a node to its successors, and
    constructing new labels for the successors, combine these new labels  
    with their previously created label lists. 
    Let node $k$ be a node with new labels. 
    %After constructing new labels for node $ k$ in the network 
    %from the extendable labels of all nodes for which the node $k$ is a successor, 
    %we combine these new labels of node $k$ with the label list of $k$ that includes previously 
    %created labels. 
    To combine the label list of $i$ with the list of new labels created 
    for $i$, compare each label with the others in the list and only keep the 
    {\em non-dominated labels}. 
    let $l^i_1$ and $l^i_2$ be two labels for node $i$.  Label $l^i_1$ 
    dominates label $l^i_2$ if all of the following conditions are 
    satisfied: 
    \begin{itemize}
    \item[i.] the current travel time and the vehicle load 
      of $l^i_1$ are less than or equal to those of $l^i_2$,
    \item[ii.] the reduced cost of $l^i_1$ is less than or 
      equal to the reduced cost of $l^i_2$, and
    \item [iii.] the number of total unreachable customer nodes 
      of $l^i_1$ are less than or equal to those of $l^i_2$, and
      for all nodes $g \in I$, $g$ is reachable only in label 
      $l^i_1$ or $g$ is unreachable in both labels $l^i_1$ and  $l^i_2$. 
    \end{itemize}
    %%%%%%%%%%%%%%%%%%%%%%%%%% COMMENT %%%%%%%%%%%%%%%%%%%%%%%%%
    \begin{comment} 
      Let $L_k^1$ and $L_k^{2}$ be two labels in the list.  Let 
      $(T_k^{1v},T_k^{1t},s_k^1,V_k^{11},..,V_k^{1N})$ and 
      $(T_k^{2v},T_k^{2t},s_k^2,V_k^{21},..,V_k^{2N})$  be the 
      corresponding state variables 
      and $C^1$ and $C^{2}$ be the reduced costs of the labels. 
      Label $L_k^1$ {\em  dominates} label
      $L_k^{2}$, if all of the following conditions are satisfied:
      \begin{itemize}
      \item The current resource consumptions of $L_k^1$ are less 
	than or equal to the resource consumptions
	of $L_k^{2}$, i.e., $T_k^{1v}\leq T_k^{2v}$ and $T_k^{1t}\leq T_k^{2t}$.
      \item The number of unreachable nodes in $L_k^1$ is less 
	than or equal to the number in $L_k^{2}$, i.e.,
	$s_k^1\leq s_k^2$.
      \item $V_k^{1p}\preccurlyeq V_k^{2p} \ \ \forall p=1..N$, 
	where the relation $V_k^{1p}\preccurlyeq V_k^{2p}$ is 
	true if $V_k^{1p}=0$ (p is reachable from k) in $L_k^{1}$ 
	or $p$ is unreachable from $k$ in both labels.
      \item The reduced cost of $L_k^1$ is less than or equal to the 
	reduced cost of $L_k^{2}$, i.e.,
	$C^1 \leq C^{2}$.
      \item The labels are not equal, i.e.,
	$(C^1,T_k^{1v},T_k^{1t},s_k^1,V_k^{11},..,V_k^{1N})$ $\neq$ 
	$(C^2,T_k^{2v},T_k^{2t},s_k^2,V_k^{21},..,V_k^{2N})$.
      \end{itemize}
    \end{comment}
    %%%%%%%%%%%%%%%%%%%%%%%%%%%%%%%%%%%%%%%%%%%%%%%%%%%%%%%%%%%%%%%%%%%5
    All dominated labels are deleted from the label lists of the nodes. 
  \item If label list is changed for any node after domination algorithm 
    is applied i.e., new labels are added, then the algorithm returns 
    to the extension step. When all the labels are extended and no new 
    labels are added to the lists of nodes, the algorithm terminates. 
  \end{enumerate}
\end{enumerate}

{\bf Additional Domination Rule}\\
In addition to \citet{Feillet}'s domination rules, we have developed
a domination rule for our pricing problem. Solving described 1p-ESPPRC on the 
constructed network leads to paths from source to sink that correspond 
to pairings including sets of {\em ordered} vehicle routes.  
Since the ordering of the routes within a pairing is arbitrary, 
we need to eliminate generating multiple paths that are
different permutations of a set of routes and correspond to same
pairing.  
%a path from source to sink in the constructed network corresponds to 
%a pairing including an ordered set of vehicle routes. 
For example, a call to 1p-ESPPRC may generate some or all of 
the following paths:
\begin{itemize}  
\item (s - 1 - 2 - si - 3 - 4 - si - 5 - 6 - si), 
\item (s - 1 - 2 - si - 5 - 6 - si - 3 - 4 - si),
\item (s - 3 - 4 - si - 1 - 2 - si - 5 - 6 - si), 
\item (s - 3 - 4 - si - 5 - 6 - si - 1 - 2 - si), 
\item (s - 5 - 6 - si - 1 - 2 - si - 3 - 4 - si), and
\item (s - 5 - 6 - si - 3 - 4 - si - 1 - 4 - si). 
\end{itemize}
However, all six 
of those paths are equal with respect to total cost and time, and they 
correspond to the same pairing composed of the following three 
routes: (facility - 1 - 2 - facility), (facility - 3 - 4 - facility) and 
(facility - 5 - 7 - facility). 
%Thus, we need to generate only one of these paths. 
%Since in the pricing problem we are 
%looking for a set routes instead of one route  - the algorithm is
%not elementary with respect to the sink node - the state space of the problem is
%larger than an ordinary ESPPRC. 
%Although the ordering of the routes within a pairing
%is arbitrary, a path from source to sink in the constructed network corresponds to 
%a pairing including an ordered set of vehicle routes. 
To eliminate generation of all possible ordered sets, we introduce a new 
domination rule in which we force a specific order for routes in the pairing, 
which decreases the state space of the algorithm significantly. We add
the following condition to the algorithm. 
%The following domination rule is added to the step 3 of the label correcting 
%algorithm.
\begin{itemize}
\item [{\bf Rule}] The index of the first customer node of each route should be less 
than the index of the first customer node in the following route.
\end{itemize}
The modified algorithm only generates paths from source to sink in which the first 
node of each route is less than the first node of the following route. 
Thus, the algorithm with this rule generates only path 
(s - 1 - 2 - si - 3 - 4 - si - 5 - 6 - si). 
%In the other paths, the first nodes of the routes are not in increasing order.  

To apply this rule, we add variable $R^i_{f}$ to keep the index of the 
first customer node in the current route in the path to the definition of
label $l_r^i$. We let  
$l_r^i = (R^i_c, R^i_{v},R^i_{t}, R^i_n, R^i_{f}, V^i_1,..,V^i_{|I|})$
be the label corresponds to path $r$ to node $i$.
We do the following changes in 1p-ESPPRC:
%Based on this definition and the domination rule we want to add, we
%add the following to the algorithm  1p-ESPPRC: 
\begin{itemize}
\item {\em In extension step of the algorithm}, in part (b), while constructing 
  label $l^k_p$ from $l^i_r$, we set $R^k_{f} = k$ if $i$ is the source or 
  sink ($k$ will be the first node in the new route), and $R^k_{f} = R^i_{f}$ (still
  in the same route, thus the first node is same) 
  otherwise. We also need to modify updating 
  vector $V^k$. 
 \begin{eqnarray*}
   V^k_g & = & \left\{ \begin{array}{llll}
             N & \mbox{ if $g = k$ and $N = \#$ of nodes in the path,}  \\
         V^i_g & \mbox{if  $V^i_g \neq 0$ and $V^i_g \neq -2$,}\\
	    -1 & \mbox{if  $V^i_g = 0$ or $V^i_g = -2$, and $R^k_{t} + T_{k,g} + T_{g,si} > L^T$,}  \\
            -2 & \mbox{if  $k$ is sink, $V^i_g = 0$ and $g < R^k_{f}$,}  \\
             0 & \mbox{otherwise,}
      \end{array} 
      \right\} \forall g \in I, \\
     R^k_n    & = &  \sum_{g \in I: V^i_g \neq 0 \mbox{ or }V^i_g \neq -2} 1. 
    \end{eqnarray*}
 Here, $V^k_g = -2$ means that node $g$ is  {\em temporarily unreachable} 
 because based on new rule, it can not be the first node in the new route.
 However, it may become the part of the
 path during further extensions. Therefore, when we calculate number
 of unreachable nodes  $R^k_n$, we do not count temporarily unreachable
 nodes, we treat them same as reachable nodes. In addition, while extending
 a path to other nodes after the first node in the route, we change the condition
 temporarily unreachable to reachable by setting $V^k_g$ to 0. 
 
%  If $k$ is the sink node, $V^k_g = 0$ for some $g \in I$ and  
%  $g < R^k_{f}$, we set  $V^k_g = -2$ meaning that the customer $g$ is
%  {\em temporarily unreachable} because of the new rule. And if $i$ is a sink 
%  node, i.e., $k$ is the first node in the route and $V^k_g = -2$
%  (set to be temporarily unreachable)  for $g \in I$, 
%  remove 
\item {\em In the elimination step of the algorithm}, the steps and the conditions that
should be satisfied to decide dominated labels are same. While executing condition 
(iii) in part (a), if there are temporarily unreachable nodes, we count them 
as reachable. Only addition to this step is when both labels dominate each other.
In the previous version, we can pick one of  $l^i_1$ or  $l^i_2$  arbitrarily 
and delete the other one. However, in this version, if  $l^i_1$ dominates 
$l^i_2$ and $l^i_2$ dominates $l^i_1$ based on conditions (i), (ii) and 
(iii), we keep the one with smallest $R^i_{f}$ (first node in the current route).  
\end{itemize}

We have updated our algorithm with this additional rule. 
%When the algorithm terminates, the labels in the list for the sink node corresponds 
%to feasible pairings. 
At the end of the algorithm, each label at the sink node corresponds to a 
feasible pairing with an associated reduced cost. The costs of paths at the sink 
node are updated with the vehicle fixed cost that must be included in the
reduced cost of a variable, and any path with negative reduced cost is a 
candidate pairing to add to the RMP. 
%The labels with negative reduced cost are candidate pairings to 
%be added to the RMP. 
If the algorithm concludes with no negative reduced cost labels on the sink node, 
then we can conclude that the current RMP is optimal. 


\subsection{2p-ESPPRC: Two-Phase Pricing Problem}
An alternative approach to solve the exact pricing problem is to extend the first 
approach (1p-ESPPRC) to a two-phase pricing problem: first generate a set of feasible 
vehicle routes and then combine them to construct a feasible pairing of minimum 
reduced cost. Each phase is also formulated as a network problem and solved as an ESPPRC.

Since a pairing is a set of routes, we can express the variable part of the 
reduced cost of a pairing as the sum of the reduced costs of the routes 
forming the pairing. For there to exist a negative reduced cost pairing 
composed of more than one route, there must exist negative reduced cost routes.
Otherwise, the pairing including only the single route with the lowest reduced
cost would be the lowest cost pairing overall.  Thus, the basic idea of the
two-phase approach is to generate feasible routes first. If there are negative reduced 
cost routes, then combine these routes to generate pairings with more negative reduced 
cost.  Otherwise, update the reduced cost of the first-phase routes with the fixed cost 
to obtain a column (a single-route pairing) with negative  reduced cost or to conclude 
that the current solution is optimal.  

\paragraph{Phase One: Generating Routes}
To generate only the feasible routes instead of pairings, we solve a regular 
ESPPRC problem on the constructed network with $|I| + 2$ nodes. We delete the
arcs from the sink to customers and we terminate the paths at sink and do
not extend them further from sink. We use algorithm 1p-ESPPRC without 
the additional domination rule. At the end of the algorithm, each label at the 
sink node corresponds to a feasible vehicle route with an associated reduced cost.  

Let $L_S$ denote the set of feasible routes with negative reduced cost at the
sink. If $L_S$ is non-empty, we continue with phase two. Otherwise, we add the
fixed cost from expression (\ref{red-cost}) to the cost of the smallest
reduced cost route to obtain the reduced cost of that single-route pairing.
If this reduced cost is zero or positive, then we can conclude that the current 
LP relaxation solution is optimal. Otherwise, we have a pairing with negative reduced 
cost to add to the RMP. Note the implication that, even in the case where there are no
\emph{routes} with negative reduced cost, we may still get a negative reduced
cost \emph{pairing} (consisting of a single route) when the fixed cost
associated with the vehicle is negative. This can occur after applying certain
branching rules (see Section \ref{branchingRule}).

\paragraph{Phase Two: Generating Pairings}
The objective of phase two is to generate feasible negative reduced cost pairings
by combining the routes in set $L_S$. Two (or more) vehicle routes can be combined to 
form a feasible pairing if they do not have customer nodes in common and they can both 
(all) be completed within the vehicle's time limit. Finding a feasible pairing with the 
minimum reduced cost again corresponds to solving a ESPPRC.

We construct a network with $2 + |L_S|$ nodes: a source node, a sink node and 
one ``intermediate'' node to represent each of the vehicle routes in set $L_S$.  
The network includes an arc from the source node to each intermediate node, arcs 
between pairs of intermediate nodes, and an arc from each intermediate node 
to the sink node.  Note that, because the sequence of routes in a pairing
is arbitrary, we only include an arc from an intermediate node to another 
intermediate node with higher index.  

Associated with each arc inbound to an intermediate node are two values: the time 
required to complete the corresponding route and the reduced cost of the route.  
The time and cost for arcs inbound to the sink are zero.  To find a feasible pairing 
with the minimum reduced cost, we solve an ESPPRC with a single resource constraint 
(time). As in phase one, we apply the extended label correcting algorithm of 
\cite{Feillet}, 1p-ESPPRC.  At the end of the algorithm, each label at the sink node corresponds 
to a feasible pairing with an associated negative reduced cost. The costs of 
pairings at the sink node are updated with the fixed cost that must be included in the 
reduced cost of a variable, and any pairing with negative reduced cost is a candidate 
pairing to add to the RMP. 


Since we solve the exact ESPPRC algorithm using 1p-ESPPRC in each phase of the  
algorithm, the overall pricing algorithm is exact and results in a column with the 
most negative reduced cost, if any. We refer to this exact pricing 
algorithm as 2p-ESPPRC.

\zaupdate{Any proposition can be stated?}

\zaupdate{Example, 2 figures representing phase 1 and phase 2}

\subsection{Bidirectional Pricing Problem}

\section {SPPRC: Solving the Relaxation of Pricing Problem}
One approach to find negative reduced cost pairings is to solve a relaxation of 
the pricing problem in which we allow customer nodes to be visited more than once,
i.e., we relax the constraint (\ref{pricing-j-1}) in the pricing problem formulation 
in Section \ref{form-pricing-prob}. Even though we relax this constraint in the pricing
problem and generate columns based on the relaxed model, we can still obtain
feasible solutions to the AUG-SPP-LRSP as long as the integrality constraints are 
satisfied. This is because the constraint (\ref{spp_con1}) and the integrality of
the variables make sure that only the columns visiting customers at most once
are selected in a feasible solution. To incorporate solutions of relaxed pricing
problem into the AUG-SPP-LRS formulation, we change the definition of the parameter 
$a_{ip}$ to represent the number of visits to customer node $i$ in the pairing $p$ 
for all $i \in I$, $p \in P_j$ and $j \in J$. Thus, the parameter $a_{ip}$ may 
take values greater than 1. If we relax other constraints in pricing problem such 
as the vehicle capacity, time limit constraint or the construction of feasible 
routes, then these columns in AUG-SPP-LRSP may lead to solutions that does not 
satisfy the constraints of an LRSP. 

Solving a pricing problem in which we allow multiple visits to customers
in the constructed network is equal to solving SPPRC. Relaxing elementary property 
for customers in ESPPRC leads to SPPRC. As noted in section (\ref{solv-pri-exact})
the complexity of a labeling algorithm for SPPRC is pseudo-polynomial, which is an
advantage over ESPPRC.  However, a disadvantage of generating columns from a
relaxed pricing problem is that, because one  set of constraints is relaxed, 
the LP relaxation bound may be weakened. That is, when the COLGEN terminates, the 
LP relaxation solution may contain columns visiting some customers multiple times, 
in which case the objective function is only a lower bound on the optimal LP 
relaxation bound obtained using an exact pricing problem. 
\zaupdate{In section \ref{comp-espprc-dp}, we demonstrate empirically how 
the bound of this relaxation compares to the optimal LP relaxation bound.} 

To solve SPPRC, we use a generic label correcting algorithm based on 
\citet{Desrochers}. The algorithm 1p-ESPPRC described in Section  (\ref{1pespprc})
can be modified to solve SPPRC. In this version, since we do not have elementary
property for customers in the path, we do not have to mark the customers 
as unreachable. Again, we let $l_r^i$ be a label corresponds to 
path $r$ from source to node $i$. We state that vector 
$l_r^i = (R^i_c, R^i_{v},R^i_{t},  V^i_1,..,V^i_{|I|})$ keeps the information
about path $r$ where $R^i_c$, $R^i_{v}$, and $R^i_{t}$ are as described in
Section  (\ref{1pespprc}). We let $V^i_k$ be the number of times customer
$k$ is visited in path $r$. Since visiting a customer is not one of the
resources in SPPRC, the algorithm only updates vector $V^i$ and does not 
compare or check it during the algorithm, thus we have the following modifications 
to algorithm 1p-ESPPRC to solve a SPPRC:
\begin{itemize}
\item While executing {\em extension step,} in part (a), the condition (i)
is deleted, and to decide whether the label $l^i_r$ for node $i$ can be extended 
to node $k$, resource feasibility (ii) for the extension is checked. 
In part (b) of this step, besides $R^i_c$, $R^i_{v}$, and $R^i_{t}$, we update 
$V^k_g$ if $g = k$ and set $V^k_g = V^k_g +1$. 
\item In {\em elimination step} in part (a) condition (iii) is deleted
since we do not compare entries of vector $V^i$ to decide domination relation
between two labels.  
\end{itemize}


%The advantage of the relaxed version is that it can be solved using a dynamic 
%programming algorithm. The complexity of the algorithm 
%is pseudopolynomial depending on the number of resource constraints (\cite{Beasley}). 
%In our case, we have two resources: the vehicle capacity ($VCap$) and the time limit 
%($TL$). Therefore, the algorithm has $O(|TL| |VCap| |N|^2)$ complexity. A disadvantage 
%of this approach is that, because one  set of constraints is relaxed, the LP relaxation 
%bound may be weakened. That is,  when the CG algorithm terminates, the LP relaxation 
%solution may contain walk columns, in which case the objective function is only 
%a lower bound on the optimal LP relaxation bound. In section \ref{comp-espprc-dp}, 
%we demonstrate empirically how the bound of this relaxation compares to the optimal 
%LP relaxation bound. 

\begin{comment}
To solve the relaxed version, we have developed a standard forward dynamic programming 
algorithm, which finds minimum cost walks from the source to the sink considering the 
vehicle capacity and the time limit constraints. \citet{Beasley} give a formal description 
of such a dynamic programming algorithm. We have three state variables: the current node, 
the accumulated time and the current vehicle load. The state variable representing the 
accumulated time increases  whenever we move from one node to another. However, the 
current vehicle load may increase or decrease depending on the node to which we arrive. 
Since we assume that the problem is a delivery problem, the vehicle load decreases when 
we arrive to a customer node.  The vehicle load increases and is set to the full truckload 
when we arrive to the sink node. To reset the vehicle load to the full truck load, we set 
the demand of the sink to $-VCap$. If we let $(i, T^t, T^v)$ be the state variables where 
$i$ represent the current node, $T^t$ represents the accumulated time from the source 
to the node $i$ and  $T^v$ represents the current vehicle load, then $c(i, T^t, T^v)$ 
is defined as the state cost. The forward recursive function is:
\begin{eqnarray*}
&c (source,0,0)    =   0 \hspace{3.02cm} & \\
&c (j, T^t, T^v)   =   \min_{i} \{ c'_{ij}+c(i,T_i^t,T_i^v)| &T_i^t+t'_{ij} =
 T^t \land T^t+t'_{j,sink}\leq TL \land \\
                        & &\min(VCap,T_i^v-d_j)=T^v \land T^v \geq 0\}
\end{eqnarray*}
The first and the third constraints in the definition of the recursive function update 
the state variables. The third constraint guarantees that the vehicle load is reduced 
when $j$ is a customer node, and the vehicle load is set to vehicle capacity if $j$ is 
the sink. The second and the last constraints are used to check the feasibility of the 
state. 
\end{comment}

\section {Solving the Pricing Problem Heuristically}
\label{heur-pri-problems}
An exact algorithm is guaranteed to generate a column with the minimum reduced cost.  
The downside of the approach is that it may not be very efficient for large problems.  
As the number of customers, vehicle capacity and time limit increase, the state space
of the algorithm, i.e., the possible number of partial paths, so the number labels in the 
label list of each node in the algorithm increases. Thus, processing all these labels 
and terminating the algorithm take a lot more time considering the complexity of an 
ESPPRC problem. 

Generally it is common in the literature to use heuristics for pricing problem to improve 
the efficiency of a branch and price algorithm. To find feasible pairings more quickly, 
here we define \zaupdate{two} heuristic versions of a labeling algorithm that we use in our 
solution algorithm. We can use these heuristic modifications for both 1p-ESPPRC
and 2p-ESPPRC.  

\subsection{ESPPRC-LL: ESPPRC with Label Limit}
This heuristic version is first proposed by \citet{Dumitrescu}
and restricts the state space of the exact algorithm by directly restricting the number 
of labels. In particular, during 1p-ESPPRC, or either phase of 2p-ESPPRC, we can limit 
the maximum number of labels at each node that are stored to be extended, i.e., number 
of  non-processed labels. During the algorithm, as we update the list of labels in 
elimination step, at the end of the domination algorithm, we order the non-dominated labels 
in {\em increasing order of reduced cost} and we keep at most the first {\em label limit} 
number of non-processed  labels. We refer to these heuristic versions of pricing algorithms 
as 1p-ESPPRC-LL(n) and 2p-ESPPRC-LL(n), where n is the value of the label limit. For small 
label limits, the algorithm terminates very quickly. And as n gets larger, the
heuristic version approaches to the exact algorithm. 

\subsection{ESPPRC-CS: ESPPRC for a Subset of Customers}
This second heuristic restricts the state space by considering just a 
subset of the customers instead of all of them. In this way, the total number of labels 
to be generated is restricted in advance. 

In each pricing iteration for a facility, we determine {\em two customers} that have 
the {\em lowest reduced costs from the facility}. Let $j \in J$ be the facility index,
we choose customers $i_1$ and $i_2$ such that 
$c^{'}_{j,i_1} \leq c^{'}_{j,i_2} \leq c^{'}_{j,i}$ for all  $i \in I \backslash \{i_1,i_2\}$. 
Let $I_n = \{i_1, i_2\}$ be the subset. We initialize the subset with these two
customers since we expect to have {\em at least two routes in a pairing}.  Then, by 
exploring the links from the subset of customers to the remaining customers, we 
repeatedly add one customer that is connected by a cheapest link to the subset until 
the size of the subset is equal to the determined {\em subset size}.
Particularly, we let $I_n = I_n \cup \{k\}$ where 
$k = \mbox{argmin}_{k \in I \backslash I_n}\{\mbox{min}_{i \in I_n}\{c^{'}_{i,k}\}\}$ 
We stop when $|I_n| = n$ where n is the determined subset size. 

After we determine $I_n$, we run exact pricing algorithm with this smaller customer
set. We refer to these heuristic versions 
of pricing algorithms as 1p-ESPPRC-CS(n) and 2p-ESPPRC-CS(n), where n is the subset size. 
In 2p-ESPPRC-CS(n), we use the selected customers in phase 1 and, state space of phase 
2 reduces depending on the reduction in the state space of phase 1.  


%\paragraph{One Pricing Problem for All Facilities}
%\zaupdate{ - Solve pricing for one facility}\\
%\zaupdate{ - Use the column generated for on facility for
%other facilities by recalculating distance} \\
 

If a heuristic version of the pricing algorithm identifies columns with negative 
reduced cost, they are feasible and they can be added to the RMP. 
However, these columns are not necessarily the columns with most negative
reduced cost for a given dual vector. Therefore, if a heuristic version of the pricing algorithm, 
do not identify any columns with negative reduced cost, the exact pricing problem
must be applied at least once in order to conclude whether or not there exist columns 
with negative reduced cost, i.e., to conclude that the current LP is optimal or not. 

\section{Branching Rules}
\label{branchingRule}

An important step in developing a branch-and-price algorithm is devising good branching rules.
Branching rules that affect the structure of the pricing problem may make column generation more
difficult. Therefore, the main objective is to find disjunctions that eliminate as much of the integer
infeasible region as possible but that minimize the increase in difficulty of solving the pricing
problem. We apply four branching rules in our branch-and-price algorithm.

%Devising good branching rules is an important step in developing a branch-and-price algorithm.
%The branching rules may affect the structure of the pricing problem, which may cause the problem 
%to become more difficult to solve than the original pricing problem. Therefore, the main 
%objective is to find branching rules that eliminate as much of the integer infeasible region 
%as possible but that do not increase the difficulty of the pricing problem too much. We apply 
%the following branching rules in our algorithm.
\begin {itemize}
\item [{\bf Rule 1}]

{\bf Branching on the facility location variables $t$:} At any node of the branch-and-price tree, 
if the corresponding LP solution has at least one fractional facility location variable, then
we apply standard variable dichotomy branching for the facility location variables.
we select one of the fractional facility location variables, say $t_j$, as our branching
variable. 
\begin {itemize}
\item In the left branch, we set $t_j=1$, which forces facility location $j$ to be open. 
  %meaning that the corresponding facility is forced to be open. 
  The pricing problem for the corresponding node does not change since $t$ variables
  are not present in the constraints of the pricing problem, and the pricing problem for 
  each facility is independent from the others. 
\item In the right branch, we set $t_j=0$, which forces facility j to be closed. 
  Since the facility is closed, no pairings associated with the facility are allowed in the 
  solution, so there is no need to solve the pricing problem. In addition, we set all $z_{p}$ 
  variables for facility $j$ to 0.
  %If facility j is closed, then no pairings associated with facility j are allowed in the 
  %solution. Thus, we fix the variables corresponding to pairings of facility j to zero. The structure 
  %of the pricing problem does not change. In fact, we do not solve the pricing problem for facility j, 
  %since it is fixed closed.
\end{itemize}
%from paper
%First, we apply standard variable dichotomy branching for the facility location variables $t$. 
%For a fractional variable $t_j$, in the right branch, we set $t_j = 1$, which forces facility location 
%$j$ to be open. The pricing problem for the right branch does not change since $t$ variables
%are not present in the constraints of the pricing problem, and the pricing problem for each facility
%is independent from the others.  In the left branch, we set   
%$t_j = 0$, which forces facility location $j$ to be closed.  Since
%the facility is closed, no pairings associated with the facility are allowed in the solution, so there
%is no need to solve the pricing problem. In addition, we set all $z_{p}$ variables for facility $j$
%to 0.

\item [{\bf Rule 2}]

  {\bf Branching on the total number of vehicles at each facility:} 
  %By definition, each pairing   corresponds to one vehicle. Also, vehicles (pairings) belong 
  %to a single facility. Thus, in 
  %any integer solution, the number of vehicles (pairings) used for any facility must be integer. 
  %If we define $nV_j$ to be the number of vehicles used at facility j, then we have the following 
  %set of equalities: 
  By definition, a pairing is associated with a single facility and corresponds to one vehicle.  
  Thus, in any integer solution, the number of vehicles (pairings) used at any facility must be 
  integer. Let $w_j$ be the number of vehicles used at facility $j$.  Then,  
  \begin{eqnarray}
    \label{n-equal}
    w_j =\sum_{p \in P_j} z_{p} \qquad \forall j \in J. 
  \end{eqnarray}
 
  If we add the equalities (\ref{n-equal}) to the master problem, we can branch explicitly on
  the $w_j$ variables. Assume that in a fractional solution $w_j = b$, where $b$ is not integer.
  We create two new nodes:
  \begin{eqnarray*}
    w_j & \leq & \lfloor b \rfloor \qquad  \mbox {(in the left branch)}\\
    w_j & \geq & \lceil b \rceil    \qquad \mbox {(in the right branch)}
  \end{eqnarray*} 
  %In one branch, we require $w_j\geq \lceil b \rceil$; in the other
  %branch, we require  $w_j\leq \lfloor b \rfloor$.
  Branching on the $w_j$ variables instead of implicitly  branching on $\sum_{p \in P_j}z_{p}$ 
  is an example of explicit constraint branching, an idea developed by \citet{Appleget} to 
  improve the performance of branch-and-bound algorithms for mixed integer linear problems. 
  These new variables make it easier to identify the necessary updates to the pricing problem. 
  Let $\beta_j$ be the dual variable associated with constraint (\ref{n-equal}) for facility $j$.
  Then, the reduced cost of the variable $z_{p}$  changes to the following:
  %At any node of the branch-and-price tree, if any of the $nV_j$ variables have fractional values, then 
  %  we select   one as our branching variable. If the selected variable $nV_j$ has the fractional value 
  %  $a$, then  we create two new branches:
  %To do this branching, we introduce the variables $nV_j$, $\forall j\in M$, and add the equations 
  %(\ref{n-equal}) to our formulation. Branching on the $nV_j$ variables instead of implicitly  branching
  %on $\sum_{p \in P_j}Z_{jp}$ is an example of explicit constraint branching, an idea developed by 
  %\citet{Appleget} to improve the performance of the branch-and-bound algorithms for mixed integer
  %linear problems. One benefit of these new variables is that it is easier to see the necessary 
  %updates to the pricing problem. Let $v_j$ be the dual variable associated with constraint 
  %(\ref{n-equal}). Then, the reduced cost of the variable $Z_{jp}$  changes to the following:
  %\begin{eqnarray}
  %  \label{new-red-cost}
  %  \hat{C}_{jp}=C_{jp}-\sum_{i\in N}a_{ipj}.\pi_{i}+\sum_{i \in N}a_{ipj}.d_{i}.\mu_{j}+\sum_{i \in N}a_{ipj}.\sigma_{ji}-v_j
  %\end{eqnarray} 
  \begin{eqnarray}
    \label{new-red-cost}
    \hat{C}_{p}=C_{p}-\sum_{i\in N}a_{ip} \pi_{i}+\sum_{i \in N}a_{ip} D_{i} \mu_{j}+\sum_{i \in N}a_{ip} \sigma_{ji}-\beta_j
  \end{eqnarray} 
  In equation (\ref{new-red-cost}), we can interpret $\beta_j$ as a fixed cost for facility $j$ 
  that is independent of pairing $p$ and the included customers.  Therefore, we do not need 
  to change our network structure or the pricing problem structure to incorporate $\beta_j$. 
  Instead, we can calculate the cost of a path in the modified network in the same way as before and 
  then add the fixed cost $-\beta_j$ to the cost of the solution to calculate the correct
  reduced cost of the column.     

  %In equation (\ref{new-red-cost}), $v_j$ can be considered as a fixed cost for facility j that is 
  %shared with every customer node in the pairing p. We do not need to change our network structure or the pricing 
  %problem structure to incorporate $v_j$. We calculate the cost of a path in the modified network in same way as 
  %before the branching rule and then add the fixed cost $-v_j$ to the cost of the shortest path 
  %to find the true reduced cost of the column with the branching rule.    

\item[{\bf Rule 3}]
  {\bf Branching on assignment of a customer to a specific facility:} 
  At any node of the branch-and-price tree, if the solution is fractional and Rules 1 and 2 are not 
  applicable, then we consider the assignment of a customer to a specific facility. 
  In a fractional solution, a customer may be partially served by two or more pairings assigned 
  to different facilities. In this case, we can create a branching rule to force a customer to 
  be assigned to a specific facility. To do so, for a customer $i$ served by more than one 
  facility and a facility $j$ such that $\sum_{p \in P_j }a_{ip} z_{p}$ is fractional, we create 
  two branches. 

  \begin{itemize}
  \item In the right branch, we force customer node $i$ to be served by facility $j$ by adding 
    the inequality
    \begin{eqnarray}
      \label{con-rul4}
      \sum_{p \in P_j}a_{ip}z_{p} \geq  1 \ \ (\gamma_{ji})
    \end{eqnarray}       
    to the RMP. Let $\gamma_{ji}$ be the dual variable associated with inequality (\ref{con-rul4}).  
    To incorporate this into the pricing problem for facility $j$, we change the cost of any arc 
    entering node $i$ to $c_{ki}-\pi_{i}+D_i \mu_j +\sigma_{ji}-\gamma_{ji}$ and solve the pricing 
    problem as before. In addition, we can remove customer $i$ from the network for the pricing
    problems of other facilities.
    
  \item In the left branch, we forbid the assignment of customer node $i$ to facility $j$ by adding
    to the RMP the constraint 
    \begin{eqnarray}
      \sum_{p \in P_j }a_{ip}z_{p} \leq  0. \nonumber 
    \end{eqnarray}
    To incorporate this restriction into the pricing problem for facility $j$, we remove node $i$ 
    from the network of facility $j$. The pricing problems for the other facilities do not change. 
   \end{itemize} 
  %\begin{eqnarray}
  %\label{branchingrule3-cond}
  %\sum_{p \in P_j }a_{ipj} Z_{jp}
  %\end{eqnarray}  
  % is fractional. We create two new branches:
  %  \begin{eqnarray*}
  %  \sum_{p \in P_j }a_{ipj} Z_{jp} & \geq & 1 \qquad  \mbox {(on the left branch)}\\
  %   \sum_{p \in P_j }a_{ipj} Z_{jp} & \leq & 0   \qquad \mbox {(on the right branch)}
  %  \end{eqnarray*}    
  
\item[{\bf Rule 4}]
  {\bf Branching on flow between pairs of customer nodes} 
  At any node of the branch-and-price tree, if the solution is fractional and Rules 1, 2  and 
  3 are not applicable, then we may encounter the case where a customer is covered by two or 
  more (fractional) pairings associated with the same facility.  Then, we branch on flows on 
  single arcs between pairs of customer nodes using a modified version of the \citet*{Ryan} 
  branching rule first suggested by \citet{Desrochers4}. 
  %then we use the modified
  %\citet*{Ryan} branching rule developed by \citet{Desrochers4}. 
  In the original \citet*{Ryan} branching rule, a pair of set partitioning constraints is 
  selected. In one branch, the rule is to branch on the sum of the variables that satisfy both 
  constraints and, in the other branch, the rule is to branch on the sum of the variables that do 
  not satisfy both constraints simultaneously. 
  %In the LRS context, the rule is to select 
  %a pair of customer nodes and to branch on the sum of pairings such that the pair of customer  
  %nodes appear in the same pairing simultaneously or not. 
  In VRP context, set partitioning constraints are associated with customers, and \citet{Desrochers4} 
  modify the Ryan and Foster branching rule by looking only at the variables in which the selected 
  pair of customers are in consecutive order or not. This modification makes it possible
  to incorporate branching rule to the pricing problem without changing the structure or
  the solution method.
  
  To apply the modified Ryan and Foster branching rule, we select a pair of customer nodes $(n_1,n_2)$ 
  such that  
  \begin{eqnarray}
    \label{branchingrule4-con}
    0 < \sum_{j \in J}\sum_{\{p \in P_j | (n_1\to n_2) \in p\}} Z_{jp} < 1
  \end{eqnarray}       
  where $\{p \in P_j | (n_1\to n_2) \in p\}$ is the set of pairings (of facility $j$)
  in which $n_1$ is followed directly by $n_2$ for some $n_1$, $n_2$ $\in I$. 
  %  Based on the structure of the pairings in the LRS problem, if $n_1$ is a customer node, then $n_2$ 
  %  can be another customer node or it can be the sink; if $n_1$ is the source, then $n_2$ can only be 
  %  a customer node and if $n_1$ is the sink, then $n_2$ can only be a customer node.  For example, if 
  %  we let $i$ and $k$ be two customer nodes, then the $(n_1,n_2)$ pair can be one of the following, 
  %  $(source,i)$, $(i,k)$, $(i,sink)$ and $(sink,i)$. 

  In one branch, we require that the flow on the directed arc between the selected pair of customers 
  is one in all possible solutions. In the other branch, we forbid the flow on the directed arc 
  between the selected pair of customers. We create two new nodes: 
  \begin{optprog}
   &\sum_{j \in J}\sum_{\{p \in P_j | (n_1\to n_2) \in p\}} Z_{jp} & \geq  & 1  & \mbox {(in the left branch)} \nonumber\\
   & \sum_{j \in J}\sum_{\{p \in P_j | (n_1\to n_2) \in p\}} Z_{jp} & \leq  & 0  & \mbox {(in the right branch)} \nonumber 
    %& \mbox{In the right branch} & \sum_{j \in M}\sum_{\{p \in P_j | (n_1\to n_2) \in p\}} Z_{jp} \geq  1 \nonumber 
  \end{optprog}       
  
  \begin{itemize}
  \item The left branch specifies that either node $n_1$ is followed by node $n_2$ in a pairing or that neither 
  node appears. To enforce this rule for the pairings to be generated, we update the arcs in the network 
  for the pricing problem. In the pricing problem, all arcs leaving node $n_1$ and all arcs entering 
  $n_2$ are deleted except the arc $(n_1,n_2)$. Therefore, if a pairing includes one of the nodes, then it 
  includes the arc $(n_1,n_2)$. To delete an arc in the network, we set the cost of the arc to infinity. 
  In the RMP, to satisfy the branching rule for the previously created pairings,
  we set to zero the variables corresponding to pairings that include only one of the nodes, 
  $n_1$ or $n_2$, or that include both nodes but in which $n_2$ does not immediately follow $n_1$.  
  
  \item The right branch specifies that $n_2$ cannot follow $n_1$ in a feasible pairing. The feasible pairings 
  may include $n_1$ or $n_2$ only or may include both of them as long as $n_2$ does not follow $n_1$ directly. 
  To apply this rule for the pairings to be generated, we simply delete the arc $(n_1,n_2)$ 
  from the network in the pricing problem. In the RMP, we set to zero the previously 
  created pairings in which $n_1$ is followed by $n_2$.
\end{itemize}
\end{itemize}

These four branching rules guarantee that the branch-and-price algorithm
terminates with an optimal solution. We further investigate 
this statement in the next section. 

\section {Terminating The Branch-and-Price Tree}
\label{BPAlg-Opt}
In a standard branch-and-bound or a branch-and-price algorithm, as in algorithm 
BRANCH-PRICE, the nodes satisfying one of the following conditions can be fathomed:
\begin{enumerate}
\item The LP for the node is proved to be infeasible.
\item The lower bound for the node is larger than the best integer feasible 
  solution. 
\item The LP solution of the node is integer. 
  %Therefore, either the solution is a candidate solution or the node was fathomed 
  %because its objective function value was larger than the best integer feasible solution.  
\end{enumerate}
In our branch-and-price algorithm, there may exist some nodes that satisfy 
none of the conditions (1) to (3), and even though the LP relaxation is 
non integer, none of the four branching rules described in Section \ref{branchingRule} 
is applicable. Therefore, these nodes satisfy the following:
\begin{itemize} 
\item [i.] the LP solution is feasible but not integer, and it is not fathomed 
  because of its objective function value, 
\item [ii.] all the facility location variables are integer (branching rule 1), 
\item [iii.] the number of vehicles for each facility is integer (branching rule 2), 
\item [iv.] each customer is assigned to a single facility (branching rule 3), and 
\item [v.]there are no pairs of customer nodes that satisfy condition  
  (\ref{branchingrule4-con}) (branching rule 4) .
\end{itemize}

We will show that the nodes satisfying these conditions can still be fathomed 
based on fathoming condition (3). We can construct an integer solution from the the LP 
solution with the same objective function value, and may update the best integer 
solution if it is a better upper bound.

In the LP solution vector satisfying conditions (ii) to (v), each open facility
is fully utilized and the set of customers is divided into disjoint non-empty 
subsets each of which is assigned to a single open facility. Therefore, the
LP solution vector can be decomposed into smaller solution vectors one for each 
facility. Each of these vectors corresponds to a solution to a single facility
LRSP instance and thus they can be investigated separately. Based on this
observation, further in this section we investigate the solution vector associated 
with a single facility. The statements are true for a solution vector corresponding 
to any facility. In this solution, the facility variable and the number of vehicles
used (total number of pairings selected) are integer, but at least two 
pairing variables are non zero and fractional. By rearranging the solution values 
and construction of fractional pairings, we can modify this solution to an integer 
solution vector associated with the same facility. To explain this, we will
show that the fractional pairings which are a part of the LP solution vector 
satisfying conditions (ii) to (v) have a certain structure.

Let $r_1$ and $r_2$ be two routes that can be included in a feasible pairing,
and $S_1$ and $S_2$ be the set of customers visited by route $r_1$ and $r_2$, 
respectively. 
\begin{defn}{\bf Node disjoint:}
%We say that 
Paths $r_1$ and $r_2$ are  node disjoint if they do not have a common customer 
node, e.g., $S_1\cap S_2 =\emptyset$.
\end{defn}
\begin{defn}{\bf Identical (Same):}
%We say that 
Paths $r_1$ and $r_2$ are  identical if they both include exactly the same set of
customers as well as same set of links between these customers. 
\end{defn}
We can represent a pairing in terms the combination of routes included in the 
pairing. The routes included in a pairing are node disjoint by construction
since in a pairing, a customer may be visited at most once. 
However, in a non-integer solution vector for a facility, at least two pairings
are fractional and thus one or more customer nodes may be included in more than 
one pairing. 
%Otherwise, the solution will be integer. 
And, in fact, at least two pairings corresponds to fractional
variables 
%for the same facility 
must include either same route(s) or different but non-node disjoint routes,
or both. We explain this result using an example. Solution 1 is an example of a 
fractional solution (for a single facility) that includes pairings that share only 
identical routes, all routes are node disjoint. Solution 2 is an example of a 
fractional solution that includes pairings that have non-node disjoint routes.   

\begin{example}
{\bf Solution 1: Same or node disjoint routes:}
%\noindent \underline{Solution 1: Same or node disjoint routes:}
\begin{eqnarray*}
  z_1 = 0.5 &: & \mbox{facility$_1$} - 1 - 2 - \mbox{facility$_1$} \\ 
  z_2 = 0.5 &: & \mbox{facility$_1$} - 1 - 2 - \mbox{facility$_1$} - 3 - 4 - 5 - \mbox{facility$_1$}\\
  z_3 = 0.5 &: & \mbox{facility$_1$} - 6 - 7 - \mbox{facility$_1$} - 3 - 4 - 5 - \mbox{facility$_1$} \\
  z_4 = 0.5 &: & \mbox{facility$_1$} - 6 - 7 - \mbox{facility$_1$}
\end{eqnarray*}
%\end{example}
%\begin{example}
\hspace{0.7in}{\bf Solution 2: Different and non-disjoint paths:}
%\underline{Solution 2: Different and non-disjoint paths:}
\begin{eqnarray*}
z_1 = 0.5 &: & \mbox{facility$_1$} - 1 - 2 - \mbox{facility$_1$} - 3 - 4 - 5 - \mbox{facility$_1$}\\
z_2 = 0.5 &: & \mbox{facility$_1$} - 4 - 6 - 7 - \mbox{facility$_1$} \\
z_3 = 0.5 &: & \mbox{facility$_1$} - 5 - \mbox{facility$_1$} - 7 - 1 - 2 - 3 - \mbox{facility$_1$} 
\end{eqnarray*}
\end{example}

Both solution 1 and 2 are fractional since at least two variables associated with
pairings including same customers are selected. However, in solution 1, 
selected pairings not only include same customers but also include same 
routes such as route facility$_1$ - 1 - 2 - facility$_1$.
%For example, both pairing 1 and 2 includes route facility$_1$ - 1 - 2 - facility$_1$.   
On the other hand, in solution 2, pairings do not include same routes
but there are at least two pairings (such as pairings 1 and 2) that
include non-node disjoint routes. 

A fractional solution for a facility can be similar to either solution 1 or
solution 2, or it can be combination of solution 1 and solution 2.
However, for a solution vector satisfying conditions (i) to (v) we
can state the following result.

\begin{lemma} 
\label{ip-lp-rel}
In a solution vector satisfies conditions (i) to (v),  
there exists at least one facility with an associated fractional solution 
vector. And a fractional solution vector for such facility corresponds to pairings 
including identical or node disjoint routes (as in solution 1). 
At least two pairings have at least one identical route. And all non identical 
routes in the pairings are node disjoint. % from each other.
\end{lemma} 

\begin{proof} 
We have noted before that a solution vector satisfying conditions (ii) to (iv), 
can be decomposed by facility. Since the solution is fractional (condition (i)), 
there exists at least one facility with a fractional solution vector. 
Since the solution vector for facility is fractional, there are at least two 
pairing variables that are fractional (location variable and the total number 
of vehicles for the facility must be integer because of conditions (ii) and (iii)).  
Because of fractional pairing variables, there exists at least one customer 
visited in two fractional pairings. 

Let $i$ be a customer included in both pairings $z_1$ and $z_2$. Let $r_1$ 
be the route visiting customer $i$ in $z_1$ and  $r_2$ be the route visiting 
customer $i$ in $z_2$. Since $r_1$ and $r_2$ are both visiting customer $i$, 
they must be either the same route or they are not node disjoint. 
If they are not identical and not node disjoint, then there exists
at least one pair of customers $k_1$ and $k_2$ such that the link $(k_1,k_2)$ 
is either in route $r_1$ or in $r_2$, but not in both. However, this contradicts 
with the satisfaction of condition (v) (branching rule 4). 
Therefore, $r_1$ and $r_2$ must include same set of links, thus they are
identical. Therefore, if two fractional pairings visit same customer, they have 
to include exactly same route(s) including common customers. This concludes
that a fractional set of pairings must have at least one route included in 
at least two pairings, and non identical routes must be node disjoint.  
\qed
\end{proof}

\begin{comment}
\noindent \underline{Solution 1: Same and disjoint routes:}
0.5 : f acility1 - 1 - 2 - f acility1
0.5 : f acility1 - 1 - 2 - f acility1 - 3 - 4 - 5 - 6 - f acility1
0.5 : f acility1 - 7 - 8 - 9 - f acility1 - 3 - 4 - 5 - 6 - f acility1
0.5 : f acility1 - 7 - 8 - 9 - f acility1
\underline{Solution 2: Different and non-disjoint paths:}
0.5 : f acility1 - 1 - 2 - f acility1 - 3 - 4 - 5 - 6 - f acility1
0.5 : f acility1 - 4 - 6 - 7 - f acility1 - 8 - 9 - f acility1
0.5 : f acility1 - 5 - 8 - f acility1 - 7 - 1 - 2 - 3 - f acility1 - 9 - f acility1

\begin{lemma} If a leaf node satisfies condition 4, then the fractional 
solution at the node includes 
%pairings that have disjoint paths. Since the solution is fractional, 
at least two pairings that have at least one identical path. And all the paths in 
the pairings are node disjoint from each other (similar to solution 1).
\end{lemma} 
\begin{proof} 
Assume that we have a leaf node satisfying condition 4 and
%and a non-integer solution. 
the fractional solution includes pairings 
that have non-disjoint paths. (The paths, $P_1$ and $P_2$, are non-disjoint
if $(P_1\cap P_2)\neq \emptyset$ and $P_1\neq P_2$). Therefore, there is at least one
arc $(k,i)$ and two paths $P_1$ and $P_2$ such that $i,k \in N$, 
$\{i\} \subseteq (P_1\cap P_2)$ and $(k,i) \in P_1$ but $(k,i) \notin P_2$
and, there is a $k^{'}$ such that $(k^{'},i) \in P_2$.
%Since $P_1$ and $P_2$ are non-disjoint, there should be a $k^{'}$ such that 
%$(k^{'},i) \in P_2$. 
In addition to that, 
$\sum_{j \in M}\sum_{p \in P_j}a_{ijp} Z_{jp}$ must be $1$. Therefore, 
because of the $(k^{'},i)$ link, $0 < \sum_{j \in M}\sum_{\{p \in P_j | (k\to i) \in p\}} Z_{jp} <1$
must be satisfied. This contradicts the assumption that the node satisfies
condition 4. \qed
\end{proof}
\end{comment}

\begin{lemma}
\label{ip-lp-rel2}
In a solution vector satisfies conditions (i) to (v), the LP relaxation value 
is equal to value of an integer solution.
\end{lemma} 
\begin{proof} 
Based on Lemma \ref{ip-lp-rel}, if the fractional solution satisfies conditions 
(i) to (v),  the structure of the solution vector for each facility is similar to 
solution 1, i.e pairings include only identical routes and node disjoint routes.
By eliminating identical routes from all pairings but one, we can obtain
a set of pairings which includes each route exactly once. This set of pairings
corresponds to an integer solution. The objective function value of the new
solution is same with the previous fractional solution since both solutions
have exactly same set of node disjoint routes and the flow on each node is 1
(each customer included in a route must be served fully by single facility) .
\qed
\end{proof}

Therefore, based on Lemma \ref{ip-lp-rel}, by setting the LP solution value equal 
to the value of of an integer solution, we can fathom a node satisfying conditions 
(i) to (v) based  on fathoming condition (3). This proves the following statement.

\begin{prop} The branching rules listed in Section \ref{branchingRule} 
are complete and guarantee that the branch-and-price algorithm provides
an optimal solution.
\end{prop} 

We conclude this section with the following example of an integer solution 
that can be obtained from solution 1 which satisfies conditions (i) to (v).    
\begin{example}
{\bf Solution 3: Integer solution from solution 1}
\begin{eqnarray*}
  z_1 = 0 &: & \mbox{facility$_1$} - 1 - 2 - \mbox{facility$_1$} \\ 
  z_2 = 1 &: & \mbox{facility$_1$} - 1 - 2 - \mbox{facility$_1$} - 3 - 4 - 5 - \mbox{facility$_1$}\\
  z_3 = 0 &: & \mbox{facility$_1$} - 6 - 7 - \mbox{facility$_1$} - 3 - 4 - 5 - \mbox{facility$_1$} \\
  z_4 = 1 &: & \mbox{facility$_1$} - 6 - 7 - \mbox{facility$_1$}
\end{eqnarray*}
\end{example}

 
%\zaupdate{After that is not finished. Do not read after this point}

\section{Overall Branch and Price Algorithm}
\label{sec:BP-overview}
%\section{Implementation Details}
%\zaupdate{Planning to add the section in lrs paper}
%\zaupdate{Add parts about artitifical variables}
\zaupdate{Add tailing off property, did some experiments about the duals values, 
early termination experiments?}

In this section, we discuss the overall branch-and-price algorithm
and some implementation details of the algorithm.  We implemented
the algorithm in MINTO 3.1 using CPLEX 9.1 as the LP solver.
MINTO (\citet{Savelsbergh2}) is a mixed integer solver that uses a 
branch-and-bound algorithm with an LP solver. MINTO has built-in applications 
such as preprocessing, variable and constraint generation and primal heuristics, 
but it also allows a user to design his/her own application code such as 
a branch-and-cut and a branch-and-price that is specific to a problem. 
Within MINTO, we customized several routines and added new ones for our algorithm.
We have written routines such as to generate an initial set of columns, to solve 
the pricing problem and to select a branching variable. In the following,
we explain some parts of the algorithm in detail. 


\subsection{Initialization}

To initialize the algorithm BRANCH-PRICE, we need to construct an initial RMP that 
includes location variables and an intial set of pairing variables. In addition,
the solution of the initial RMP  must be feasible to be able to call COLGEN.  
In this section, we describe the implementations in order to satisfy these
requirements. 

\subsubsection{Artificial variables} 
At every node of the branch-and-price tree including the root node, we need to start 
the column generation algorithm with a feasible RMP or we need to know that the 
node is infeasible to be able to fathom it. Since the RMP does not contain all of the columns,
it is not necessarily true that the real LP at the  node is infeasible when the RMP is 
infeasible. At the root node, it is easy to start with a feasible LP solution, but the 
RMP might become infeasible with respect to the current set of columns as the branching
rules are added. At a non root node, depending on the branching rules applied, the 
infeasibility may not be easily detected by checking simple conditions. Therefore,
it may not be obvious to move a feasible solution or conclude the 
infeasiblity of the current LP. If the current RMP solution at any node is infeasible, 
we need to verify that it truly is infeasible so we can fathom the node or we need to find a 
feasible solution from which to start the column generation algorithm to continue processing
the node.    


% only to the facility location variables, 
%it is easy to detect the infeasibility with respect to the number of facilities. By 
%considering fixed and unfixed location variables,  we can simply check whether the 
%lower bound on the number of selected facilities (constraint (\ref{lb-open-fac})) 
%is satisfied or not.
%One reason is that the number of pairing variables, 
%the summation of which determines the number of vehicles at each facility, is large.
%Another reason is that to decide whether there is a feasible solution with given upper bounds and
%lower bounds for the number of vehicles in facilities, one must check the set partitioning constraints
%(\ref{spp_con1}) as well as the facility capacity constraints (\ref{spp_con2}). 


We introduce artificial variables to model SPP-AUG-LRSP to be able to have an initial 
feasible RMP or obvious infeasibility for any node. We can add the artificial variables 
at the root node or at any node in which the RMP is infeasible. Currently, we introduce them 
at the root node and keep them along the tree. 
%With artificial variables, we are able to find an initial feasible solution at any node for 
%any set of restrictions on the facility variables, the total number of vehicle variables and 
%the pairing variables.    

We create an imaginary facility with $I$ pairings, one for each customer. Each pairing (column) 
visits  one customer. Let $t_a$ be equal to 1 if the artificial facility is open, and 0 otherwise. 
Let $z_{a1},..,z_{a|I|)}$ be the artificial pairing variables where $z_{ak}$ represents 
the pairing visiting customer $k$. We modify the LP relaxation of AUG-SPP-LRSP with these 
artificial variables:
\begin{optprog}
Minimize & \objective{\sum_{j \in J} C^F_j t_{j} + \sum_{j \in J} \sum_{p\in P_j} C_p z_p + C^F_a t_a + \sum_{i \in I}C_{a}z_{ai}} \label{spp-obj-artupd}\\
subject to & \sum_{j \in J} \sum_{p \in P_{j}} a_{ip} z_{p} + z_{ai}& = & 1 & \forall i \in I, \label{spp_con1-artupd} \\
           & \sum_{p \in P_j} \sum_{i \in I} a_{ip} D_i z_{p} - L^F_j t_j & \leq & 0 & \forall j \in J, \label{spp_con2-artupd} \\
           & \sum_{j \in J} t_{j} & \geq & R \label{spp_con3-artupd} \\ 
           & \sum_{p \in P_{j}}a_{ip}z_{p} - t_{j} &\leq & 0 \label{spp_con4-artupd} \\ 
           & z_{ai} - t_{a} & \leq & 0 & \forall i \in I \label{spp_con5-artupd}\\
           & 0 \leq t_{j}  & \leq & 1 & \forall  j \in J, \label{binary1-artupd} \\
           & 0 \leq z_{p} &  \leq & 1 & \forall p \in P_j, j \in J, \label{binary2-artupd} \\
           & 0 \leq z_{ai} &  \leq & 1 & \forall i \in I, \\
           & 0 \leq t_a  &  \leq & 1,
\end{optprog} 
where  $C^F_a$and $C_{a}$ are costs of artificial variables. We do not actually need to 
add variable for artificial facility and constraints (\ref{spp_con5-artupd}), we do that
for the interpretation, and consistency with the model variables. We assign large costs
to the artificial variables to discourage them being non zero in the solution, unless 
it is infeasible. We define $C^F_a = (R+1) \cdot \mbox{max}_{j \in J}\{C^F_j\}$, and
$C_{a} = 2C^O\sum_{i \in I}t_{1,i} + C^V$ where 
%They eliminate any infeasibility 
%caused by the constraint (\ref{spp_con1-upd}) and/or  the branching constraints of 
%(\ref{low-bd}) and (\ref{up-bd}).  
%\begin{align*}
% FC_A & =nREQ \cdot max_{j \in M}\{FC_j\}+1 \nonumber \\
% C_A & = \sum_{i \in N}(2 \cdot t_{facility_j,i}) C^O + C^V  \nonumber
%\end {align*} 
the fixed cost of the artificial facility is set equal to the minimum number of required facilities 
plus one times the maximum fixed cost of the facilities to get a larger cost than selecting
real facilities. The total operating cost of each artificial 
pairing is the same with  a value equal to the sum of the operating costs of all single 
customer routes from depot $j$ plus the vehicle fixed cost. We can select
depot $j$ arbitrarily. 
 
The model with artificial variables results either feasible or infeasible LP relaxation solution 
for all possible sets of branching rules. If it is infeasible, we can conclude that the 
node is infeasible and can be fathomed. If it is feasible, we can start doing column
generation. And at the end of algorithm COLGEN, if the solution includes 
artificial variables, this means that the column generation algorithm searched all possible 
columns but could not find a column to replace the artificial variables, thus
the real LP is infeasible. In this case, because of the artificial variable costs, the RMP 
solution value is high enough to fathom the node based on the incumbent solution value. 


\subsubsection{Initial Columns} 
\label{sec:initcols}
Artificial variables added to the model can be used to initialize the algorithm.
From previous computational experience with branch-and-price algorithms, 
we expect a reduction in the number of pricing iterations required to
solve the LP relaxation when starting from a good set of columns rather than an easily
constructed or arbitrary solution. Therefore, besides the artificial variables, 
we generate feasible columns to construct the initial RMP. However, there is a trade off 
between the effort (time) required to generate such a solution and the resulting reduction 
in overall running time. 

We have developed an algorithm that combines a facility location heuristic, 
several well known VRP heuristics (Clark and Wright, nearest neighbor, 
nearest insertion and sweep heuristics), and a bin-packing heuristic (best fit 
decreasing) to generate an initial set of pairings. The basic idea of the routine 
is to sequentially determine the following decisions for a given set of open 
facilities: (i) allocation of customers to facilities, (ii) design of feasible 
vehicle routes, and (iii) assignment of routes to vehicles. The total cost of such 
a sequentially determined solution for a given set of open facilities is calculated. 
To decide what set of open facilities would yield such a solution of lowest cost, 
the algorithm starts with all facilities open and iteratively closes the facility 
whose elimination results in the greatest cost savings until the problem becomes 
infeasible (this is called the DROP heuristic in the context of facility location problems).
We refer to this heuristic algorithm as I-Heur. I-Heur produces a set of columns
as well as a feasible solution for the problem. The outline of the algorithm
I-Heur which is the combination of functions DROP($J$) and TCOST($J$) is given below.

{\bf I-Heur}
\begin{itemize}
\item {\bf DROP($J$)}
\begin{enumerate}
\item Let $J_o = J$ be the set of open facilities.
\item For all facilities in set $J_o$:
  \begin{itemize}
  \item[i.]Pick $k \in J_o$ and close facility $k$.
    Let $J^k = J_o \backslash \{k\}$.
  \item[ii.] Call function TCOST($J^k$) to calculate the cost of opening all
    of the facilities in set $J^k$. 
  \end{itemize}
\item Let $k^* = \mbox{argmin}_{k \in J_o}\{\mbox{TCOST}(J^k)\}$.
\item If TCOST($J^{k^*}$)$< \infty$, then 
  let  $J_o = J_o \backslash \{k^*\}$ and return to step 2.
  STOP otherwise.
\item Return set of open facilities, $J_o$, total cost of the system
  TCOST($J_o$) and set of pairings constructed in function TCOST($J_o$).
\end{enumerate}
\item {\bf TCOST($J_t$)}
\begin{enumerate}
\item Partition customers among set of facilities, $J_t$.
  Assign each customer to the nearest facility, if the capacity
  of the facility is not exceeded with the assignment. Otherwise,
  assign the customer to the next nearest facility. Let $I_j$
  be the set of customers assigned to facility $j$ $\in J_t$.
\item Let COST $= \sum_{j \in J_t}C^F$ (add the fixed cost of facilities). 
\item For each facility $j$ in set $J_t$:
  \begin{enumerate}
  \item[i.]Call function ClarkWright($j,I_j$) for VRP instance with 
    customers $I_j$  and facility $j$. This function constructs the vehicle routes
    based on Clark and Wright heuristic. Let $R_1$ be the set of routes
    generated. Call BestFitD($R_1$) to pack the routes into pairings 
    considering time limit (best fit decreasing heuristic for bin packing
    problems). Let $P_1$ be the set of pairings generated.
    Let $C_1$ be the fixed and operating costs of pairings in set $P_1$. 
  \item [ii.]Call function NN($j,I_j$) for VRP instance with 
    customers $I_j$  and facility $j$. This function constructs the vehicle routes
    based on nearest neighbor heuristic. Let $R_2$ be the set of routes
    generated. Call BestFitD($R_2$). Let $P_2$ be the set of pairings generated
    and $C_2$ be the costs.
  \item[iii.]Call function NI($j,I_j$) for VRP instance with 
    customers $I_j$  and facility $j$. This function constructs the vehicle routes
    based on nearest insertion heuristic. Let $R_3$ be the set of routes
    generated. Call BestFitD($R_3$). Let $P_3$ be the set of pairings generated
    and $C_3$ be the costs.
  \item [iv.]Call function SW($j,I_j$) for VRP instance with 
    customers $I_j$  and facility $j$. This function constructs the vehicle routes
    based on sweep heuristic. Let $R_4$ be the set of routes
    generated. Call BestFitD($R_4$). Let $P_4$ be the set of pairings generated
    and $C_4$ be the costs.
  \item[v.] Let COST = COST + min$\{C_1, C_2, C_3, C_4\}$. Let
    $P_j = P_{argmin\{C_1, C_2, C_3, C_4\}}$.
  \end{enumerate}
\item Return COST, total system cost, set of pairings $P_j$ for 
  all $j \in J_t$.
\end{enumerate}
\end{itemize}

I-Heur returns a set of columns for facilities that are selected to be open. 
We add these columns to the RMP and initialize the incumbent solution
be the cost of the solution obtained from I-Heur. To balance the number of columns 
added for each facility, for the facilities that are fixed to be closed in the heuristic 
solution, we determine a set of customers that can be served from the facility within 
the time limit and use function ClarkWright and BestFitD to determine a set of
pairings for those facilities. 

%%%%%%%%%%%%%%%%%%%%%%%%%%%%%%%%%%%%%%%%%%%%%%%%%%%%%%%%%%%%%%%%%%%%%%%%%%%%%%%%%%%%%%%%%%
\begin{comment}
We use DROP heuristic to determine a set of facilities 
to be open. Basically, the DROP heuristic starts with all facilities are being open. 
At each iteration, one facility which results the most cost saving is closed until
the problem becomes infeasible. Total cost which includes the facility fixed 
costs, vehicle routing costs and vehicle fixed costs is calculated at every 
iteration of DROP heuristic for a given set of open facilities. 

In the total cost calculation procedure (TCOST), the first step is for any open set of facilities 
to determine an assignment of customers based on the distance and facility capacity. 
For each open facility and given customer assignment, using a VRP 
heuristic vehicle routes are constructed to serve the demands of customers assigned to
the facility. After designing the routes, a bin packing heuristic (best fit decreasing)
is used to combine the routes considering the time limit for the vehicles. This
decides the assignment of vehicles to routes as well as the number of vehicles, 
therefore the fixed cost of vehicles. The total cost calculation is repeated for each 
of the following VRP heuristics: Clark and Wright, nearest neighbor, nearest insertion 
and sweep heuristic. The minimum total cost obtained from the total cost calculation
procedure using one of the VRP heuristics is returned to the DROP heuristic to 
determine the facility to be closed. 

When the DROP heuristic terminates, it
returns a set of open facilities and the total cost of the solution which includes 
the routing and scheduling information as well as the facility fixed costs.
We add the columns which represent the solution from the combined heuristic
to the RMP and update the upper bound. To balance the number of columns added 
for each facility, for the facilities that are fixed to be closed in the heuristic 
solution, we determine a set of customers that can be served from the facility within 
the time limit and use VRP heuristics and the bin packing heuristic in combination to 
determine feasible pairings. 
\end{comment}
%%%%%%%%%%%%%%%%%%%%%%%%%%%%%%%%%%%%%%%%%%%%%%%%%%%%%%%%%%%%%%%%%%%%%%%%%%%%%%%%%%%%%%%%%%

\begin{comment}
{\bf Pricing Problem}
As discussed in Section~\ref{sec:PricingProb}, solving the pricing problem exactly
by generating and processing all the labels can be time-consuming.  To reduce
the time spent, we use the two previously described heuristic versions of the pricing algorithm 
before employing the exact pricing algorithm.  
The basic steps of the overall pricing method are as follows:
\begin{enumerate}
\item Determine a subset of customers for a given subset size, n. Execute 
2p-ESPPRC-CS(n) for the selected subset. 
\item Repeat the previous step for a limited number of iterations or until the algorithm fails 
(the algorithm {\it fails} if it does not generate any negative reduced cost columns). 
\item Execute 2p-ESPPRC-LL(m), starting with a small label limit value. 
\item Increase the label limit value m when the algorithm 
fails.
\item After subsequent failures, execute 2p-ESPPRC.   
\end{enumerate} 
The steps of the pricing problem and the parameters used differ slightly in the 
different versions of the overall branch-and-price algorithm as we explain below. 
\end{comment}

\subsection{Branch Selection}
%Within the LRSP, there is a natural hierarchy of the decisions, locate, then route and then
%assign.  This hierarchy is applied to the branching as well.
Given a fractional solution at a node of the branch-and-price tree, we apply the branching
rules in the order that they are described in Section \ref{branchingRule}. For each 
branching rule, we choose the most fractional disjunction. 

\subsection{Primal Feasible Solution}
\label{sec:primal-feas}
Having good primal feasible solutions can reduce the size of the explored tree.  
We used four ways to obtain primal feasible solutions in our solution 
algorithms: MINTO's built-in primal heuristic routines, I-Heur,
a standard branch-and-bound algorithm, and a heuristic branch-and-price algorithm.

For relatively small problems, since the quality of the solution 
generated by the I-Heur is sufficiently high, we initialized the upper bound with the 
solution obtained from I-Heur and used MINTO's primal heuristic routines within the 
branch-and-price algorithm.
 
For larger, more difficult problems, we implemented a more powerful procedure to 
both improve the upper bound and generate an initial set of columns. 
We run a heuristic branch-and-price algorithm called H-BP, in which we construct
an initial branch-and-price tree subject to a time limit, employing only the
 heuristic pricing algorithms 2p-ESPPRC-LL(n) and 2p-ESPPRC-CS(n).  
When the time limit is reached (or the algorithm terminates), we have a set of 
generated columns available to be used and an upper bound for the problem.  
To improve the quality of the upper bound obtained by H-BP, we applied
a standard branch-and-bound algorithm to an integer RMP with the current
set of columns at the root node of the H-BP subject to a one-hour 
time limit. Our computational experiments showed that the quality of the 
upper bound obtained using H-BP is very high. 

%We used four ways to obtain primal feasible solutions in our solution 
%algorithm. 
%We used  MINTO's built-in primal heuristic routines within the branch-and-price
%algorithm and initialized the upper bound with the solution obtained 
%from I-Heur. For relatively small problems, the quality of the initial solution 
%generated by the I-Heur is sufficiently high. For larger, more difficult problems, 
%however, we have implemented a more powerful procedure to both improve the upper bound
%and generate an initial set of columns. 
%We run a heuristic branch-and-price algorithm called H-BP, in which we construct
%an initial branch-and-price tree subject to a time limit, employing only the
% heuristic pricing algorithms 2p-ESPPRC-LL(n) and 2p-ESPPRC-CS(n).  
%When the time limit is reached (or the algorithm terminates), we have a set of 
%generated columns available to be used and an upper bound for the problem.  
%Our computational experiments showed that the quality of the 
%upper bound obtained using H-BP is very high. 

%Another way of obtaining a feasible solution is to 
%apply a standard branch-and-bound algorithm and solve integer RMP with current
%set of columns. We apply this method only at the root node of the H-BP subject to a one-hour 
%time limit.  

\subsection{Solving Pricing Problem and Algorithm Overview}
\label{sec:alg-view}
We have implemented two versions of a branch-and-price algorithm. 
For small or medium size instances, we ran a {\em one-stage}
solution algorithm that is an exact branch-and-price algorithm 
initialized with columns generated by I-Heur. We refer to
this version as E-BP. For larger instances,
we implemented a {\em two-stage} solution algorithm designed
on the idea of ``restart''. 
The first stage of the algorithm runs H-BP to get a good upper bound and
a good set of columns. The second stage is an exact branch-and-price algorithm 
that is initialized with the results of H-BP. We refer to this version 
as E-BP-2S. In our experiments, we used E-BP for instances with 25 customers and 
E-BP-2S for instances with 40 customers.

Since the steps of the overall pricing algorithm and the parameters used in E-BP and 
E-BP-2S are different, we separately explain the pricing algorithm in each version.

\subsubsection{Pricing Algorithm In E-BP.}
\label{pri-E-BP}
 At each node of the branch-and-price tree, we employed the following procedure:
\begin{enumerate}
\item Set both the subset size and the iteration limit to 12. Execute 2p-ESPPRC-CS(12) for 
12 iterations or until the algorithm fails (the algorithm {\it fails} if it does not 
generate any negative reduced cost columns).
\item Set the starting label limit value (m) to 50 and the maximum label limit
to 300.  
\item Call 2p-ESPPRC-LL(m) for each facility. If, during an iteration, no negative reduced 
cost columns are generated for a particular facility at its current label limit, 
execute the following procedure:
\begin{list}{\labelitemi}{\leftmargin=4em} 
\item[i.] If the current label limit is less than the maximum, increment the label
limit with a multiplier of 2, i.e., let $m = 2*m$, and solve 2p-ESPPRC. 
%\item [ii.] 
If the pricing problem identifies negative reduced cost columns, keep the label 
limit value, return to the beginning of step 3 and continue with the pricing problem for 
the next facility.  Otherwise, set the label limit to the maximum. 
\item [ii.] If the current label limit is equal to or greater than the maximum, do not solve 
the pricing problem for this facility during the subsequent pricing iterations for the 
current node. Return to step 3. 
\end{list}
\item  When no columns are generated for any facility, the pricing problem for all facilities 
must be solved using the exact version, 2p-ESPPRC, to prove optimality.
\end{enumerate}

\subsubsection{Pricing Algorithm In E-BP-2S.}
In E-BP-2S, the algorithm first calls H-BP for a limited time. 
The steps of the pricing algorithm in H-BP are similar to E-BP except that some parameters 
are different and the exact pricing 2p-ESPPRC is never used in H-BP.  The details follow. 
\begin{enumerate}
\item Set the subset size to 15 and the iteration limit to 20. 
Call 2p-ESPPRC-CS(15) for 20 iterations or until the algorithm fails. 
\item Set the starting label limit value (m) to 5 and the maximum label limit
to be 15.  
\item Call 2p-ESPPRC-LL(m) for each facility. If, during an iteration, no negative reduced 
cost columns are generated for a particular facility at its current label limit, 
do the following:
\begin{list}{\labelitemi}{\leftmargin=4em}   
\item[i.] If the current label limit is less than the maximum, increment the label
limit with a multiplier of 2, i.e., let $m = 2*m$, and re-solve 2p-ESPPRC-LL(m). 
%\item[ii.] 
If the pricing problem identifies negative reduced cost columns, keep the label 
limit value, return to the beginning of step 3 and continue with the pricing problem for 
the next facility.  Otherwise, return to step i. 
\item [ii.]If the current label limit is equal to or greater than the maximum, skip  
the pricing problem for this facility during the subsequent pricing iterations for the 
current node. Return to step 3. 
\end{list}
\end{enumerate}

After solving H-BP, we initialize E-BP-2S and start the second stage of E-BP-2S in which
the pricing algorithm is the same as E-BP except that it skips step 1, i.e., it does not call 
2p-ESPPRC-CS(n). 

